\makeatletter
\def\input@path{{styles/}}
\makeatother

\documentclass[
  a4paper,
  fontsize=10pt,
  twoside=true,
  numbers=noenddot,
  fontmethod=tex
]{kaobook}

\usepackage[english]{babel}
\usepackage{csquotes}
\usepackage{graphicx}
\usepackage[maxbibnames=10,minbibnames=3]{kaobiblio}
\usepackage{makecell}
% These newtx extension glyphs otherwise copy as Latin letters (e.g. X).
\pdfglyphtounicode{summationtext}{2211}
\pdfglyphtounicode{summationdisplay}{2211}
\pdfglyphtounicode{producttext}{220F}
\pdfglyphtounicode{productdisplay}{220F}
% Emit copyable word spaces only in text. Enabling this inside math causes
% spurious glyphs with newpxmath on pdfTeX 1.40.29 (TeX Live 2026).
\AtBeginDocument{%
  \pdfglyphtounicode{space}{0020}%
  \pdfinterwordspaceon
  \everymath\expandafter{\the\everymath\pdfinterwordspaceoff\aftergroup\pdfinterwordspaceon}%
  \BeforeBeginEnvironment{equation*}{\pdfinterwordspaceoff}%
  \AfterEndEnvironment{equation*}{\pdfinterwordspaceon}%
  \BeforeBeginEnvironment{align*}{\pdfinterwordspaceoff}%
  \AfterEndEnvironment{align*}{\pdfinterwordspaceon}%
}

% algorithm2e prints a semicolon at the end of every pseudocode line by
% default. \; stays as the line terminator in the source, but nothing is
% printed, so no semicolon appears anywhere in the document.
\DontPrintSemicolon

\usepackage{needspace}
% Note: the kaobook class (styles/kao.sty) already loads algorithm2e, so the
% algorithm environment is available without an extra package here.

% Compatibility for the kaobook caption hook with current LaTeX.
\LetLtxMacro{\kaooriginalcaption}{\caption}
\NewDocumentCommand{\kaocaptionwrapper}{s o m}{%
  \IfBooleanTF{#1}{%
    \kaooriginalcaption*{#3}%
  }{%
    \IfValueTF{#2}{%
      \kaooriginalcaption[#2]{#3}%
    }{%
      \kaooriginalcaption{#3}%
    }%
  }%
  \toggletrue{kaocaption}%
}
\AtBeginEnvironment{figure}{\LetLtxMacro{\caption}{\kaocaptionwrapper}}
\AtBeginEnvironment{table}{\LetLtxMacro{\caption}{\kaocaptionwrapper}}

% Captions are typeset in the margin column, which is too narrow to justify
% cleanly: a long hyphenated compound after the bold label leaves gaps that
% TeX has to stretch. The kaobook caption formats end in \@margin@par, which
% issues \justifying and so overrides the caption package's justification
% option; the formats are therefore redeclared here with \RaggedRight after
% it. This overrides styles/kao.sty without editing it.
\makeatletter
\DeclareCaptionFormat{margin}{\@margin@par\RaggedRight #1#2#3}
\DeclareCaptionFormat{llap}{%
  \begin{lrbox}{\mycap}%
  \begin{minipage}{\marginparwidth}%
  \@margin@par\RaggedRight #1#2#3%
  \end{minipage}%
  \end{lrbox}%
  \marginnote[0.2cm]{\usebox\mycap}%
}
\makeatother

% Metric table generation and automated result reporting
% \input{includes/metric_tables.tex}  % Disabled: pgfplotstable macros need fixing

\addbibresource{main.bib}
\graphicspath{{figures/}{logos/}}
\newcommand{\myfooter}{}

% upstream tightening is still required for those.
\widowpenalty=10000
\clubpenalty=10000
\displaywidowpenalty=10000
% Let TeX stretch a little more before declaring a line overfull, which clears
% the few remaining small overfull boxes without visibly loosening the text.
\setlength{\emergencystretch}{3em}

\setcounter{tocdepth}{1}

\makenoidxglossaries
% Definition order follows the workflow: errors, sampling, grouping, scoring,
% and evaluation. The glossary is printed with sort=def in main.tex.
\newglossaryentry{llm}{
  name={LLM},
  description={large language model: a neural network trained on large text
    collections to predict tokens and generate language}
}

\newglossaryentry{hallucination}{
  name={hallucination},
  description={a broad term for an output that is unsupported, factually wrong,
    or unfaithful to a source. The evaluation here labels correctness without
    separately identifying the cause of an error}
}

\newglossaryentry{confabulation}{
  name={confabulation},
  description={a wrong answer whose meaning changes across repeated samples
    of the same prompt. Such arbitrary errors are the target of semantic entropy}
}

\newglossaryentry{temperature}{
  name={temperature},
  description={a sampling parameter that rescales token logits before they
    become probabilities. Lower values concentrate probability on likely
    tokens; higher values flatten the distribution}
}

\newglossaryentry{nli}{
  name={NLI},
  description={Natural Language Inference: classifying whether a premise
    entails a hypothesis, contradicts it, or supports neither decision (neutral)}
}

\newglossaryentry{entailment}{
  name={entailment},
  description={a relation in which one statement implies another. The
    short-answer clustering rule requires entailment in both directions,
    with the question included as context}
}

\newglossaryentry{deberta}{
  name={DeBERTa},
  description={Decoding-enhanced BERT with disentangled attention: a
    Transformer architecture. The NLI-trained checkpoints used here compare
    answer pairs for entailment}
}

\newglossaryentry{semanticcluster}{
  name={semantic cluster},
  description={a group of sampled answers treated as sharing a meaning.
    The implemented anchor scan uses mutual entailment; its assignments
    depend on answer order and do not enforce transitivity}
}

\newglossaryentry{semanticentropy}{
  name={semantic entropy},
  description={Shannon entropy over semantic clusters of repeated answers.
    The discrete version uses sample counts; the full version uses normalized
    cluster likelihood weights. A high value indicates varied meanings,
    not necessarily incorrectness}
}

\newglossaryentry{surfaceentropy}{
  name={surface entropy},
  description={entropy over counts of identical normalized answer strings.
    It groups by text rather than meaning. Counting raw strings without
    normalization gives \emph{naive sample entropy}}
}

\newglossaryentry{naiveentropy}{
  name={naive predictive entropy},
  description={also called naive entropy: the mean negative length-normalized
    log-likelihood of sampled answers. Unlike naive sample entropy, it uses likelihoods
    rather than string counts}
}

\newglossaryentry{ptrue}{
  name={P(True)},
  description={model self-evaluation of an answer using sampled alternatives
    as context. The uncertainty score is one minus the probability of the
    ``true'' option, without renormalizing over the two options}
}

\newglossaryentry{sep}{
  name={SEP},
  description={Semantic Entropy Probe: a linear classifier on fixed model
    hidden states, trained to predict high semantic entropy. After training,
    it scores one answer without generating alternatives}
}

\newglossaryentry{accuracyprobe}{
  name={accuracy probe},
  description={a linear classifier on fixed model hidden states, trained on
    judged correctness rather than semantic entropy; also called
    \emph{embedding regression}}
}

\newglossaryentry{correctnessrule}{
  name={correctness rule},
  description={the procedure assigning correct or incorrect labels to answers,
    such as reference-answer token overlap or an LLM judge. It is separate
    from the uncertainty estimator}
}

\newglossaryentry{auroc}{
  name={AUROC},
  description={area under the receiver operating characteristic curve:
    the probability that an incorrect answer has higher uncertainty than a
    correct one, counting ties as one half. A value of 0.5 is chance ranking}
}

\newglossaryentry{hierarchicalbootstrap}{
  name={hierarchical bootstrap},
  description={resampling at multiple levels. The main paired comparison
    resamples dataset--generator units, then questions within their seed
    cells, preserving the pairing between estimators}
}


\titlehead{%
\begin{center}
  \includegraphics[height=1.4cm]{KU_Logo.png}
  \quad
  \includegraphics[height=1.6cm]{MIDS_Logo_en.pdf}
\end{center}
}

% The examination office confirmed that only the English title is registered.
\newcommand{\englishtitle}{Detecting Hallucinations in Large Language Models Using Semantic Entropy}

\subject{Bachelorarbeit / Bachelor's Thesis}
\title[\englishtitle]{\englishtitle}
\author[Jakub Wiśniewski]{%
\textsc{Jakub Wiśniewski}\\
Matrikelnummer (matriculation number): [removed]\\
Studiengang (degree program): Data Science (Bachelor), 6.\ Fachsemester\\
Geburtsdatum (date of birth): [removed]
}
\date{\today}

\publishers{%
\textsc{Betreuer (supervisor): Prof.\ Dr.\ Felix Voigtlaender}\\[0.4em]
\textsc{Mathematisch-Geographische Fakultät}\\
\textsc{Mathematisches Institut für Maschinelles Lernen und Data Science}\\
\textsc{Katholische Universität Eichstätt-Ingolstadt}\\[0.6em]
\textsc{Prüfer (examiner): Prof.\ Dr.\ Felix Voigtlaender}
}

\begin{document}

\frontmatter
\maketitle

\chapter*{Abstract}
\addcontentsline{toc}{chapter}{Abstract}

Large language models generate fluent text even when their answers are
semantically inconsistent or factually unsupported, and they rarely
signal when an answer should not be trusted. This thesis studies
\emph{semantic entropy} as a training-free method for hallucination
detection and selective abstention in short-answer question answering.

Semantic entropy measures uncertainty over meanings rather than over
surface strings. Sampled answers are grouped into semantic equivalence
classes using bidirectional Natural Language Inference entailment, and
Shannon entropy is computed over cluster frequencies. The method is
designed to detect \emph{confabulations}---answers that are both
incorrect and semantically inconsistent across alternative generations
from the same prompt---while making no claim about high-confidence
systematic errors that fall outside this failure mode.

We implement discrete semantic entropy on two structurally distinct
tasks: TriviaQA, an open-domain factual question-answering benchmark,
and SVAMP, a short-answer arithmetic reasoning dataset. Experiments are
organized in two phases on NHR@FAU high-performance computing resources.
Phase~1 uses Mistral-7B with a DeBERTa NLI backend to validate the
end-to-end pipeline at modest scale. Phase~2 extends the evaluation to
Llama-3.1-70B and stronger semantic judges on A100 and H100 hardware.
Uncertainty quality is measured with the area under the receiver
operating characteristic curve (AUROC), the area under the
rejection-accuracy curve (AURAC), and rejection-accuracy curves that
compare discrete semantic entropy against surface-form entropy baselines.

% TODO: replace the paragraph below with empirical findings once
% Phase 1 and Phase 2 HPC results are available.
Preliminary infrastructure validation confirms that the pipeline
produces schema-valid output records across both datasets and that
correctness labeling and entropy scoring are end-to-end reproducible.
Full empirical results comparing semantic entropy against surface-form
baselines, together with a manual audit of semantic clustering behavior,
are reported in Chapter~6 following completion of the HPC runs.

We additionally scope two Phase~2 extensions---Semantic Entropy Probes,
which approximate semantic entropy from hidden model states, and Kernel
Language Entropy, which replaces hard clustering with a kernel-based
entropy objective---as cost-versus-quality comparators once the discrete
baseline is established. The implementation, JSONL output schema, and
evaluation scripts are made available as reproducible artefacts for
future work on semantic uncertainty in large language models.

\chapter*{Acknowledgements}
\addcontentsline{toc}{chapter}{Acknowledgements}

I would like to thank my supervisor, Prof.\ Dr.\ Felix Voigtländer of the
Mathematical Institute for Machine Learning and Data Science (MIDS) at
Katholische Universität Eichstätt-Ingolstadt, for his guidance throughout
this project, for suggesting semantic entropy as a tractable and
well-motivated research problem, and for his feedback on the thesis draft.

I am grateful to NHR@FAU for providing the high-performance computing
resources used for the Phase~1 and Phase~2 experiments on the Alex
cluster, together with the Fritz cluster for posthoc reclustering and
Kernel Language Entropy runs.

\vspace{2em}
% TODO: confirm the exact wording with the KU Eichst\"att-Ingolstadt
% examination regulations before submission and adapt if necessary.
\section*{Declaration of independent work}
\addcontentsline{toc}{section}{Declaration of independent work}

I hereby declare that I have written this thesis independently and have
not used any sources or aids other than those cited. All passages taken
verbatim or in substance from other works are marked as such, and the
sources are listed in the bibliography. I further declare that
artificial-intelligence tools were used only as auxiliary aids; in
particular, AI assistants were used for proof-reading, code
auto-completion, and language polishing of passages written and verified
by the author. No AI tool was used to generate ideas, experimental
results, or substantive content without subsequent verification. The
thesis has not previously been submitted, in this or in a similar form,
to any other examination authority.

\vspace{3em}
\noindent
\rule{6cm}{0.4pt}\hfill\rule{6cm}{0.4pt}\\
Eichstätt, \textit{[date]} \hfill Jakub Wiśniewski

\chapter*{List of aids used}
\addcontentsline{toc}{chapter}{List of aids used}

\begingroup
\RedeclareSectionCommand[beforeskip=1.2\baselineskip,afterskip=0.4\baselineskip]{section}

\section*{Literature access and reference verification}

arXiv, ACL Anthology, OpenReview and publisher websites were used to access
research papers and check publication details. Individual papers are cited
where their content is used and listed in the bibliography.

\section*{Published method and code}

The method reimplemented here is that of Farquhar et
al.~\cite{semantic_entropy_nature_2024}. The reference implementation is the
authors' released code~\cite{kossen_semantic_uncertainty_code_2024}, which was
used as the specification for clustering, prompting, P(True) scoring and
question selection. Deviations from it are recorded in
Section~\ref{sec:deviation-register}.
The reference implementation was accessed through GitHub; its archived
release on Zenodo is identified by the accompanying code citation.

\section*{Datasets}

TriviaQA~\cite{joshi_triviaqa_2017}, NQ-Open~\cite{lee_nq_open_2019} and
SVAMP~\cite{patel_svamp_2021}, prepared using a reimplementation of the
released loading and selection procedure. The biography prompts follow the
FActScore subject format~\cite{min_factscore_2023}.
Hugging Face Hub was used to obtain model checkpoints and datasets and to
consult model cards.

\section*{Models}

\begin{itemize}
  \raggedright
  \item \emph{Answer generation}, at the pinned revisions given in
    Table~\ref{tab:generators}: \texttt{Mistral-7B-Instruct-v0.3},
    \texttt{Llama-3.1-8B-Instruct} and \texttt{Llama-3.1-70B-Instruct}.
  \item \emph{Entailment}: \texttt{microsoft/deberta-v2-xlarge-mnli},
    \texttt{cross-encoder/nli-deberta-v3-large},
    \texttt{cross-encoder/nli-deberta-v3-base} and
    \texttt{Qwen/Qwen2.5-72B-Instruct}.
  \item \emph{Correctness judging}: \texttt{Qwen/Qwen2.5-72B-Instruct} and
    \texttt{Llama-3.1-70B-Instruct}.
\end{itemize}

\section*{Software}

Python with \texttt{PyTorch}, \texttt{Transformers}, \texttt{vLLM},
\texttt{NumPy}, \texttt{pandas}, \texttt{scikit-learn}, \texttt{Matplotlib},
\texttt{datasets} and \texttt{pytest}. Environments were managed with
\texttt{uv}, revision control with \texttt{Git}. The thesis was typeset in
\LaTeX\ using the \texttt{kaobook} class, with \texttt{Biber} and
\texttt{biblatex} for the bibliography. Builds were managed with
\texttt{latexmk}.

\section*{Computing resources}

The NHR@FAU Alex GPU cluster for answer generation and language-model judging,
and the NHR@FAU Fritz cluster for post-hoc cross-encoder clustering, scheduled
with \texttt{Slurm}.

\section*{Artificial-intelligence assistants}

Claude Code (Anthropic) and Codex (OpenAI) were both used as coding assistants
in the accompanying repository and to generate the figures and tables in this
thesis from the recorded measurements. As stated in the declaration, these
tools were not used
as a source of claims, results or citations. Every reported number was checked
against the recorded measurements.
\par
\endgroup

\tableofcontents
\glsaddall
\printnoidxglossary[title=Glossary, nonumberlist]

\mainmatter

\chapter{Introduction}

Large language models can answer open-ended questions fluently even when
the answer they produce is unsupported or simply wrong. Such fluent
mistakes are commonly called \emph{hallucinations}, and in the
following we use this term as an umbrella for several distinct failure
modes (misinformation, confabulation, systematic reasoning errors).
A more careful taxonomy is given in
Chapter~\ref{ch:background}; for now the term refers to any answer the
model produces that is not supported by the input or by the facts.
The practical problem this raises is not merely that an answer is wrong,
but that a useful system should also be able to indicate \emph{when} its
answer is unreliable so that a downstream user, or an automated
verification step, can act on it.

Empirical evidence confirms the scale of the risk. Under adversarial
prompting, hallucination rates in clinical decision-support settings have
been measured as high as
50--82\,\%~\cite{omar_adversarial_hallucination_clinical_2025}; this is an
adversarial upper bound rather than a typical deployment rate, but
factuality failures remain common enough in ordinary use that the broader
landscape of deployed
LLMs~\cite{augenstein_factuality_challenges_2024} presents substantial
challenges for automated verification. The precise rate matters less here
than the structural problem it creates.
The central technical difficulty is that free-form natural language has
many valid surface forms. Two
sampled answers may use very different words while expressing the same
underlying claim, and two similar-looking answers may make incompatible
claims.

This thesis is closely based on the work of Farquhar et
al.~\cite{semantic_entropy_nature_2024}, who introduced \emph{semantic
entropy} as a training-free signal for hallucination risk. The basic
idea is taken directly from that paper: instead of measuring how much
the strings produced by the model vary, one tries to measure how much
the underlying \emph{meanings} of those strings vary. If repeated samples
from the model mostly express the same meaning, the model may be
uncertain only about phrasing; if the samples split across genuinely
different meanings, the model is semantically uncertain and the answer
should be treated as riskier, following the semantic uncertainty
framing of Kuhn et al.~\cite{kuhn_semantic_uncertainty_2023} and
Farquhar et al.~\cite{semantic_entropy_nature_2024}. The
goal of the present thesis is not to invent this idea, but to
re-implement it under a smaller, locally reproducible setup and to study
how reliably it works in that setting.

\section{Research Aim}

The main research question is whether semantic entropy provides a
better signal for hallucination risk and selective abstention than
surface-form uncertainty, and how that comparison changes when the
output shifts from short answers to paragraph-length generation. The
study covers TriviaQA, NQ-Open, and SVAMP as short-answer benchmarks
and a FactScore-style biography task as the long-form regime ---
evaluated both whole-paragraph and at the level of individual claims
--- across three open generators and three decoding seeds. It therefore
targets \emph{confabulation}-like failures, where the model's answer
varies semantically across samples drawn under the same prompt, while
also testing the harder case where surface variation is nearly
universal. It does not claim that semantic entropy detects every type
of hallucination, and in particular it cannot flag high-confidence
systematic errors where the model is consistently wrong, a failure
mode discussed by
Kalai et al.~\cite{kalai_why_language_models_hallucinate_2025} and
Simhi et al.~\cite{simhi_choke_2025}.

The working hypothesis from Farquhar et
al.~\cite{semantic_entropy_nature_2024} is that discrete semantic
entropy computed over bidirectional
entailment-based clusters achieves a higher AUROC --- a ranking metric
for placing incorrect answers above correct ones --- than surface-form
entropy computed over sampled answer strings on the same prompts. A
secondary hypothesis is that the improvement, if present, is at least
directionally consistent across the short-answer benchmarks, despite
their very different answer formats (open-domain factual recall vs.\
short numeric arithmetic). The biography task tests the complementary hypothesis
that meaning-aware uncertainty methods separate more clearly once raw
surface diversity ceases to be informative.

\section{Study Structure}

The study separates cheap surface-form baselines from the semantic
clustering condition that most directly follows Farquhar et
al.~\cite{semantic_entropy_nature_2024}.

\textbf{Surface-form baselines.} For each prompt, the
generation model is sampled several times, the answers are lightly
normalized, and two answers are grouped together only if their
normalized strings are identical. We call this \emph{exact-match}
clustering. A deliberately naive companion score counts raw strings
before normalization. These baselines make the full pipeline
auditable end to end. It does \emph{not} implement true semantic
entropy in Farquhar et al.'s sense~\cite{semantic_entropy_nature_2024}, because
exact-match grouping ignores paraphrases; throughout the thesis the
exact-match number is therefore reported as a surface-form baseline rather
than as semantic entropy.

\textbf{Replication condition.} The main semantic-entropy condition repeats
the same model--dataset grid on the NHR@FAU
high-performance computing cluster and replaces exact-match clustering
with a bidirectional Natural Language Inference (NLI) clustering step
using a DeBERTa-v3-base cross-encoder. This is the configuration that
corresponds most directly to the method of Farquhar et al.

\textbf{Additional analyses.} Separate sections report judge sensitivity
checks with DeBERTa-v3-large and Qwen2.5-72B-Instruct, a continuous
kernel-based variant called Kernel Language
Entropy~\cite{nikitin_kernel_language_entropy_2024}, and held-out
Semantic Entropy Probes. These analyses extend the replication question;
they are not part of the baseline.

\section{Contributions}

The substantive method comes from Farquhar et
al.~\cite{semantic_entropy_nature_2024}; the contribution of this thesis
is to re-implement it under a smaller, fully local setup and to use that
setup to ask a question the original aggregate evaluation does not
isolate --- how the value of meaning-level clustering scales as outputs
lengthen. This work therefore sets out to deliver, mainly through
implementation, evaluation, and scope choices:

\begin{enumerate}
  \item an implementation of discrete semantic entropy for short-answer
    and long-form hallucination detection, covering answer normalization,
    exact-match clustering, bidirectional NLI clustering with DeBERTa-v3,
    entropy scoring, and correctness labeling;
  \item an experimental design that separates surface-form baselines,
    the NLI-clustered replication condition, and additional analyses,
    with Slurm templates and reproducibility logging for each run;
  \item an auditable JSONL output schema recording sampled answers,
    semantic clusters, cluster representatives, entropy scores, and
    correctness labels for every evaluated prompt;
  \item an evaluation framework for selective prediction, using AUROC,
    rejection-accuracy curves, and their area summary (AURAC),
    applicable to any uncertainty estimator that writes a score field to
    the shared schema;
  \item extension baselines and sensitivity checks covering Semantic
    Entropy Probes (hidden-state classifiers), Kernel Language Entropy,
    P(True), in- and out-of-distribution embedding regression, and
    semantic-judge comparisons across DeBERTa-v3-base, DeBERTa-v3-large,
    DeBERTa-v2-xlarge-MNLI, and Qwen2.5-72B-Instruct;
  \item a controlled short-versus-long-form comparison that places the
    FactScore-style biography task alongside the short-answer
    benchmarks --- including a claim-level long-form extension --- and
    exposes the regime where meaning-aware methods separate cleanly from
    surface-form baselines.
\end{enumerate}

\section{Thesis Structure}

Chapter~\ref{ch:background} reviews the different forms of
LLM hallucination, the role of uncertainty estimation, the notion of
semantic equivalence, and related semantic uncertainty methods.
Chapter~\ref{ch:semantic_entropy} defines the discrete semantic entropy
method on which the empirical work is based.
Chapter~\ref{ch:experimental_design} turns the method into a reproducible
experimental design, and Chapter~\ref{ch:hpc} documents the
implementation and HPC execution. Chapter~\ref{ch:extensions}
describes Kernel Language Entropy as a completed posthoc check and
Semantic Entropy Probes as a held-out probe extension.
Chapter~\ref{ch:results} reports the empirical results,
Chapter~\ref{ch:discussion} discusses interpretation and limitations,
and Chapter~\ref{ch:conclusion} states the final conclusion.

\chapter{Background}
\label{ch:background}

Large language models routinely produce fluent answers that are not
supported by the facts of the task. Such failures are typically grouped
under the broad label of ``hallucination''. The term covers several
distinct mechanisms: unfaithfulness to source material, unsupported
factual claims, gaps in the model's knowledge, faulty reasoning, and
answers encouraged by benchmarks that reward guessing over
admitting ignorance~\cite{ji_hallucination_survey_2023,%
huang_llm_hallucination_survey_2023,%
maynez_faithfulness_factuality_2020,%
kalai_why_language_models_hallucinate_2025,%
augenstein_factuality_challenges_2024}. This thesis uses the word as a
convenient cover term, but does not try to detect every kind. The target is
narrower: answers that are unreliable because the model is uncertain about
what it just said.

\section{Hallucinations and confabulations}

The semantic entropy paper~\cite{semantic_entropy_nature_2024}, on which
this thesis is based, sharpens the target to \emph{confabulations}:
fluent answers that are both incorrect and arbitrary, where ``arbitrary''
means that the answer is sensitive to irrelevant variation such as the
random seed used during decoding. In addition to confabulations, several
other types of hallucinations can occur: a model may be consistently
wrong because its training data contains a misconception, because the
problem requires reasoning it cannot perform, or because evaluation
incentives discourage abstention. Semantic entropy does not address
these cases; it specifically detects instability across sampled
meanings.

This narrower framing motivates the central hypothesis of the thesis,
again taken from Farquhar et
al.~\cite{semantic_entropy_nature_2024}: if the model knows
the answer, stochastic samples should vary mostly in wording while
preserving meaning. If the model is confabulating, the samples should
split across several incompatible meanings. Measuring uncertainty over
meanings rather than over surface strings should therefore be a more
informative signal for free-form question answering, following
semantic uncertainty and semantic entropy
work~\cite{kuhn_semantic_uncertainty_2023,semantic_entropy_nature_2024}.
In particular, uncertainty-based detection is, by construction, unable
to flag high-certainty hallucinations or systematic
errors~\cite{simhi_choke_2025,ma_semantic_energy_2025}.

\section{Uncertainty in free-form generation}

When a model has to pick from a fixed, short list of answers, say
``cat'', ``dog'', or ``horse'', measuring its uncertainty is easy: it
assigns a percentage to each option, and those percentages already say
how sure it is. A language model has no such tidy list. It instead
produces whole sentences, and the catch is that the \emph{same fact} can
be written in many ways (``Paris'', ``Paris, France'', ``the capital is
Paris''), while two sentences that look almost identical can mean
different things. So counting how often the exact words differ is a poor
guide to whether the model is really unsure. Split the model's probability
mass across ten paraphrases of one fact and the token-level numbers look
uncertain, even though the model has said the same thing ten times. Sequence-level entropy or
token-level variation can therefore be misleading for natural-language
answers, as discussed in work on uncertainty in natural language
generation~\cite{baan_uncertainty_nlg_2023} and semantic
uncertainty~\cite{kuhn_semantic_uncertainty_2023}.
The thesis distinguishes three levels of variation accordingly:

\begin{enumerate}
  \item \emph{lexical variation}: different token sequences or
    paraphrases (e.g.\ ``Paris'' vs.\ ``Paris, France'');
  \item \emph{semantic variation}: different meanings or propositions
    (e.g.\ ``Paris'' vs.\ ``Lyon'');
  \item \emph{factual correctness}: whether the chosen meaning is
    correct for the task.
\end{enumerate}

Semantic entropy operates at the second level. It estimates uncertainty
over semantic answers, not over truth directly; correctness evaluation
is a separate step that compares the selected answer against a reference
or a judge in semantic entropy~\cite{semantic_entropy_nature_2024}.
Recent surveys~\cite{shorinwa_uq_llm_survey_2024,liu_uq_confidence_calibration_survey_2025}
place
this family of methods within uncertainty quantification for LLMs,
alongside calibration, verbalized confidence, internal-state methods,
and other sampling-based estimators, and systematic benchmarks now
compare many such estimators head to
head~\cite{lm_polygraph_vashurin_2025}.

\section{Confidence, calibration, and abstention}

The motivation for estimating uncertainty is not only to attach a score
to an answer, but to decide what should happen next: answer, abstain,
defer to retrieval, request human review, or invoke a more expensive
verification procedure. A related notion is \emph{calibration}. A model is well calibrated when its
stated confidence matches how often it is actually right: among all answers
given with 90\,\% confidence, roughly 90\,\% should be correct. Calibration
and accuracy are different properties, and a model can have one without the
other. The view taken here needs neither. Selective prediction only requires
that risky answers be \emph{ranked} above safe ones. This
makes selective prediction and
rejection-based evaluation the natural evaluation lens for the thesis.
Accuracy alone is insufficient, because many benchmarks reward a lucky
guess more than an honest admission of uncertainty, a problem discussed
by Kalai et al.~\cite{kalai_why_language_models_hallucinate_2025}.
Rejection-accuracy curves and their AURAC summary are better aligned
with the practical question of whether uncertain answers can be
filtered out. Related work has also
examined whether semantic entropy can be used as a training signal to
encourage abstention: Tjandra et
al.~\cite{tjandra_abstain_semantic_entropy_2024} fine-tune models on it
directly. That is a different use of the same quantity, and it is not what
this thesis does. Here the score is only ever computed on a finished answer
from a model that is left untouched; no gradient is taken through it.

A number of related works have approached the same problem without
explicit semantic clustering. P(True)~\cite{kadavath_language_models_know_2022}
asks the model to judge whether
its own answer is correct, probing verbalized self-knowledge rather than
sample diversity. SelfCheckGPT~\cite{manakul_selfcheckgpt_2023}
instead samples multiple responses and checks them for consistency
without relying on external databases or access to token
probabilities. Semantic entropy is
closest in spirit to these sampling-based methods, but, unlike them, it
explicitly aggregates samples into semantic clusters before computing
an entropy score~\cite{kuhn_semantic_uncertainty_2023,semantic_entropy_nature_2024}.
These methods span a spectrum, from purely internal (P(True) reads the
model's own verdict) to purely behavioral (SelfCheckGPT compares samples).
Semantic entropy sits between them. It adds one step none of the others
perform --- deciding when two answers mean the same thing --- and the value
of exactly that step is what this study measures.

A third family reads the model's \emph{hidden states} rather than its words.
The observation behind it is that a model's internal activations often carry
information about the reliability of an answer that its output text does not
express. Burns et al.~\cite{burns_latent_knowledge_2023} recover a
truth-like direction in activation space without any labels; Azaria and
Mitchell~\cite{azaria_internal_state_lying_2023} train a small classifier on
hidden states to tell true from false statements; and
INSIDE~\cite{chen_inside_2024} scores hallucinations from the geometry of
sentence embeddings taken from inside the model. The practical appeal is cost:
these methods need one forward pass, not $M$ samples and $M^2$ entailment
comparisons. Semantic Entropy Probes~\cite{kossen_semantic_entropy_probes_2024}
join the two families by training exactly such a probe to predict semantic
entropy itself, so that the expensive signal is paid for once at training time
and read cheaply thereafter. Both a probe of this kind and a
correctness-trained variant are evaluated in
Chapter~\ref{ch:results} as single-pass comparators; they are defined in
Sections~\ref{sec:sep-method} and~\ref{sec:other-estimators}. Note that
training such a probe is a separate, supervised step layered on top of a frozen
generator, and it does not change the generator either.

\section{Semantic equivalence and entailment}

For computing the semantic entropy, the key operational step is
deciding when two sampled answers should be treated as the same answer.
Semantic entropy adopts semantic equivalence rather than exact string
equality: in Farquhar et al.'s original
formulation~\cite{semantic_entropy_nature_2024}, two answers belong to
the same group when they entail one another in the context of the question.
This bidirectional-entailment
criterion is a practical approximation to equality of meaning.

Natural Language Inference (NLI) is the standard tool for this. An NLI
model reads two sentences and labels their relationship as
\emph{entailment} (the first implies the second), \emph{contradiction},
or \emph{neutral}. Checking entailment in both directions is then a
practical stand-in for ``these two answers mean the same thing.'' The
specific NLI models used here are DeBERTa-v3 \emph{cross-encoders} (in
two sizes, base and large) fine-tuned on the standard SNLI and MultiNLI inference
datasets~\cite{sentence_transformers_nli_deberta_v3_base_2022,%
he_deberta_2021,he_debertav3_2021,williams_multinli_2018}. ``Cross-encoder''
just means the two sentences are fed into the model
\emph{together}, so it can directly compare them; the alternative, a
\emph{bi-encoder}, would encode each sentence on its own and compare the
results afterwards. The cross-encoder is more accurate for this kind of
pairwise judgement but slower, because every pair of answers has to be
run through the model.
The exact bidirectional-entailment rule and the recorded backend choices
are settled in Chapters~\ref{ch:semantic_entropy} and~\ref{ch:results}.

\section{Research gap}

Prior evaluations of semantic entropy report strong aggregate AUROC but
leave three questions unresolved. First, the original evaluation compares
semantic entropy against naive entropy over raw sampled strings, but not
against a normalized surface-form
baseline~\cite{semantic_entropy_nature_2024}, and the survey and
benchmark literature that follows it inherits that set of
comparisons~\cite{shorinwa_uq_llm_survey_2024,%
liu_uq_confidence_calibration_survey_2025,lm_polygraph_vashurin_2025}. It
is therefore unclear how much of the
reported gain over ``naive'' entropy is attributable to semantic
clustering rather than to simple answer normalization. Second, short-answer and long-form results are usually reported together
rather than as a controlled contrast. Yet these are exactly the two regimes
in which surface-level and meaning-level uncertainty should come apart.
Recent work does report that semantic entropy weakens as one-sentence outputs
grow longer~\cite{nguyen_beyond_semantic_entropy_2025}, but it does not
isolate the short-versus-long contrast under one protocol. Third,
the dependence of sampling-based detection on the decoding temperature,
the very parameter that controls sample diversity, has been
examined only in isolated ablations (including in the same
work~\cite{nguyen_beyond_semantic_entropy_2025}) rather than
characterized systematically across tasks and detector families. This
thesis targets these three gaps directly.

Addressing them fixes the empirical scope to three short-answer tasks,
namely open-domain factual question answering with
TriviaQA~\cite{joshi_triviaqa_2017} and NQ-Open~\cite{lee_nq_open_2019}
and arithmetic word problems with SVAMP~\cite{patel_svamp_2021}, and
one long-form task, a FactScore-style biography
task~\cite{min_factscore_2023}. The dataset rationale, the datasets left
outside the recorded study, representative records, and the exact prompt
construction are given in Chapter~\ref{ch:experimental_design} and
Appendix~\ref{app:prompts}.

Taken together, the literature places semantic entropy between cheap
surface signals and expensive factual verification. Testing it fairly means
keeping generation, clustering, scoring and evaluation separate, so this
thesis builds a small, verifiable implementation first.

\chapter{Semantic Entropy and Related Estimators}
\label{ch:semantic_entropy}

This chapter defines the method: discrete semantic entropy over sampled
answers. The aim is to keep the central idea of Farquhar et
al.~\cite{semantic_entropy_nature_2024}, which is to measure uncertainty over
meanings rather than over strings, while stating the procedure precisely enough that
the numbers in Chapter~\ref{ch:results} can be traced to specific steps of it.
The meaning-level framing follows Kuhn et
al.~\cite{kuhn_semantic_uncertainty_2023}.

Two commitments run through the chapter. Where the published description and
the released code~\cite{kossen_semantic_uncertainty_code_2024} disagree, the
code is what gets described, because a reader reproducing the result runs the
code and not the prose. The clustering step in Section~\ref{sec:clustering} is
where that matters. And the pipeline is kept in separable pieces, namely sampling,
grouping and scoring, so Chapter~\ref{ch:experimental_design} can vary one
while holding the others fixed. The last sections
(\ref{sec:related-estimators}--\ref{sec:other-estimators}) define the three
estimators compared against it.

\section{Entropy as a measure of uncertainty}
\label{sec:entropy-measure}

For a discrete distribution $p_1,\ldots,p_K$ over $K$ outcomes, the Shannon
entropy is
\[
  H(p) \;=\; -\sum_{k=1}^{K} p_k \log p_k .
\]
It is small when almost all the mass sits on one outcome and large when the
mass spreads over many. The maximum is $\log K$, at the uniform distribution.
The base of the logarithm changes the unit but not the ordering, so it does not
affect AUROC. This thesis uses the natural logarithm, so values are in nats.
Entropy is used because it is the quantity this literature defines and
reports~\cite{kuhn_semantic_uncertainty_2023,semantic_entropy_nature_2024}.
Nothing turns on the choice: AUROC depends only on the ranking, so any
increasing measure of spread would give similar results.

For semantic entropy the same formula applies, but the outcomes are not raw
strings. They are groups of strings taken to mean the same thing, so the
entropy measures how spread out the meanings are.

\subsection{Notation}

Table~\ref{tab:notation} lists the symbols used in this chapter.

\begin{table}[ht]
\thisfloatsetup{capposition=above}
\centering
\caption{Notation used for discrete semantic entropy.}
\label{tab:notation}
\begin{tabular}{@{}lp{8.2cm}@{}}
\hline
\textbf{Symbol} & \textbf{Meaning} \\
\hline
$x$ & prompt (question) supplied to the model \\
$M$ & number of sampled answers drawn from the model per prompt \\
$y_i$ & the $i$-th sampled answer, $i \in \{1,\ldots,M\}$ \\
$Y(x)$ & the multiset of all $M$ sampled answers for prompt $x$ \\
$\ell_i$ & length-normalized sequence log-likelihood of $y_i$ \\
$e(y_i, y_j)$ & NLI entailment decision: 1 if the entailment model labels
                        the pair $y_i \to y_j$ as \emph{entailment},
                        0 if it labels it \emph{neutral} or
                        \emph{contradiction} \\
$\sim_{\mathrm{sem}}$ & semantic equivalence relation used by the
                        clustering step (see Section~\ref{sec:clustering}) \\
$C(x)$ & the set of clusters $\{c_1,\ldots,c_K\}$ induced on $Y(x)$ \\
$K$ & number of distinct clusters, $1 \le K \le M$ \\
$n_k$ & number of sampled answers assigned to cluster $c_k$ \\
$\hat{p}_k$ & empirical cluster probability $n_k / M$ \\
$\widehat{H}_{\mathrm{sem}}(x)$ & discrete semantic entropy score \\
$\widehat{H}_{\mathrm{lex}}(x)$ & surface entropy score over normalized
                                  strings (the clustering-free comparator;
                                  called \emph{surface entropy} in
                                  Chapter~\ref{ch:results}) \\
AUROC & area under the receiver operating characteristic curve \\
AURAC & area under the rejection-accuracy curve \\
\hline
\end{tabular}
\end{table}

\section{The method, step by step}
\label{sec:method-steps}

The method has three computational steps. All three run on a fixed, trained
model and only read what it produces; none of them changes its weights.

\begin{enumerate}
  \item \textbf{Sampling.} For a prompt $x$, draw $M$ answers
        $y_1,\ldots,y_M$ under a fixed decoding configuration. The values used
        are in Section~\ref{sec:sampling-settings}.
  \item \textbf{Clustering.} Group the $M$ answers into $K$ clusters, putting
        two answers together if they are judged equivalent in the context
        of $x$.
  \item \textbf{Entropy scoring.} Compute each cluster's share
        $\hat{p}_k=n_k/M$ and the entropy
        $\widehat{H}_{\mathrm{sem}}(x)=-\sum_k \hat{p}_k\log\hat{p}_k$.
\end{enumerate}

Attaching a correctness label and using the score to decide what to answer
belongs to the \emph{evaluation}, not to the method, and
Chapter~\ref{ch:experimental_design} handles it separately. Keeping that line
is what makes the comparisons in Chapter~\ref{ch:results} possible: clustering
consumes stored samples, so the entailment model can be swapped without
regenerating anything, and correctness labelling consumes a finished answer, so
the grading rule can be swapped without touching the estimators.

Figure~\ref{fig:method-schematic} reproduces the original illustration. Panel
(a) makes the point on a short question: counting raw strings treats ``Paris''
and ``It's Paris'' as different answers and the entropy looks too high, while
grouping by meaning gives a low score. Panel (b) shows the long-form version,
where a biography is split into claims, each claim is re-asked as a question,
and the regenerated answers are grouped by meaning.

\begin{figure}[ht]
\centering
\includegraphics[width=\textwidth]{semantic_entropy_paper_fig1.png}
\caption{The semantic entropy method. (a)~On a short question, count\-ing
  distinct wordings makes the model look more un\-cer\-tain than it is; the
  source figure labels this ``mis\-lead\-ing\-ly high naive entropy''.
  Grouping by meaning first gives a low semantic entropy. (b)~The same idea
  applied to paragraph-length bi\-og\-ra\-phies, one claim at a time.
  Reproduced from Farquhar et al.~\cite{semantic_entropy_nature_2024}
  (\emph{Nature}, 2024), licensed under CC~BY~4.0.}
\label{fig:method-schematic}
\end{figure}

\section{Sampling setup}
\label{sec:sampling-setup}

For each prompt the model is sampled $M$ times under one decoding
configuration, giving $Y(x) = \{y_1,\ldots,y_M\}$. Two parameters control the
randomness. \emph{Temperature} $T$ divides the token logits before the softmax.
Low $T$ gives little variety, high $T$ gives more. \emph{Nucleus sampling}
(top-$p$) keeps only the smallest set of tokens whose probability adds up to at
least $p$.

Because the estimators read nothing but these samples, the temperature decides
how much signal there is to measure at all.
Section~\ref{sec:temperature-sensitivity} reports what happens across the
range, and Section~\ref{sec:sampling-settings} gives the values used.

\section{Clustering by equivalence}
\label{sec:clustering}

The main design choice is the equivalence relation. Two are used.

\paragraph{Normalized exact match.} The simplest relation:
$y_i \sim_{\mathrm{exact}} y_j$ when their normalized strings are identical.
Normalization is each dataset's standard answer-cleaning step. For the
question-answering datasets it lowercases, strips punctuation and the articles
\emph{a\,/\,an\,/\,the}, and collapses whitespace, following the TriviaQA
evaluation protocol, so ``The Eiffel Tower!'' and ``eiffel tower'' become the
same string. For the arithmetic dataset it extracts the final number, so ``The
answer is 7 apples'' becomes ``7''. This is a surface grouping: it collapses
spelling and formatting and nothing else. The entropy over these groups is
called \emph{surface entropy}, and it is the clustering-free comparator that
RQ2 turns on. A cruder companion, \emph{naive sample entropy}, counts the raw
strings without normalizing first.

\paragraph{Bidirectional entailment.} Instead of string identity, this rule
asks whether two answers entail each other. For $y_i$ and $y_j$ it asks
whether $y_i$ entails $y_j$ \emph{and} whether $y_j$ entails $y_i$. If both
hold, they are the same semantic answer:
\[
  y_i \sim_{\mathrm{sem}} y_j
  \quad\Longleftrightarrow\quad
  e(y_i,y_j)=1 \ \text{and}\ e(y_j,y_i)=1.
\]
Farquhar et al.~\cite{semantic_entropy_nature_2024} define equivalence
\emph{in the context of the question}, and give the reason: ``Paris'' does not
entail ``The capital of France is Paris'' on its own, but in the context of the
question it does. The short-answer results here follow that definition. Every
entailment query in the canonical configuration is prefixed with the question,
and the LLM backend is given the question alongside both answers for the same
reason. The one exception is the whole-paragraph biography protocol, which is
answer-only, because a generated paragraph is not a response to a single
question.

Two details of the input are easy to assume wrongly. The entailment step
consumes the \emph{raw} generations, not the normalized strings above:
normalization belongs to the comparator, and applying it first would confound
the two. And the question is joined to each answer as plain text, so what the
entailment model reads is a pair of question-prefixed strings.

The rule approximates equivalence rather than defining it, because judgments
from a small cross-encoder can be noisy. It may split answers that differ only
in harmless detail, and merge answers that hide real differences. Four backends
compute $e(\cdot,\cdot)$:

\begin{itemize}
  \item \texttt{microsoft/deberta-v2-xlarge-mnli}, the model the released code
    loads by default~\cite{he_deberta_2021,williams_multinli_2018}. This is the
    \emph{canonical} backend: results reported without a named backend use it,
    because it is what a reader running the published code would get.
  \item \texttt{cross-encoder/nli-deberta-v3-large}, a current
    cross-encoder~\cite{sentence_transformers_nli_deberta_v3_large_2022,%
    he_debertav3_2021}.
  \item \texttt{cross-encoder/nli-deberta-v3-base}, a smaller
    cross-encoder~\cite{sentence_transformers_nli_deberta_v3_base_2022}, used
    on the long-form arm only.
  \item \texttt{Qwen/Qwen2.5-72B-Instruct}~\cite{yang_qwen25_2024}, an open
    stand-in for the proprietary judge Farquhar et al.\ use on sentence-length
    answers. It is prompted with the entailment instruction taken verbatim from
    the released code.
\end{itemize}

Every backend is applied afterwards to one stored set of generations, so the
entailment model is the only thing that differs between them and any change in
the scores comes from it alone. Section~\ref{sec:clustering-sensitivity}
quantifies how much that choice matters, and it turns out to matter more than
the difference between the estimators.

\subsection{The grouping procedure as implemented}
\label{sec:anchor-scan}

The grouping follows the released
code~\cite{kossen_semantic_uncertainty_code_2024} rather than a tidier
equivalent. Answers are scanned in generation order. The first answer without a
cluster identifier becomes an \emph{anchor} and gets a new identifier, and
every later answer mutually entailed with that anchor gets the same one. The
scan then moves to the next unassigned answer.

Two properties matter, and neither is incidental.

First, the assignment in the inner loop is unconditional. The released routine
guards only the \emph{outer} loop, which decides whether an answer becomes an
anchor; the inner assignment carries no such test. An answer a previous anchor
already claimed is therefore \emph{reassigned} when a later anchor also matches
it, so its final cluster is decided by the last anchor that matches rather than
the first.

Second, the procedure computes no transitive closure. If A is mutually entailed
with B, and B with C, but A and C are not, then A, B and C need not end up in
one group; whether they do depends on which the scan reaches first. The result
is \emph{order-dependent}: permuting the samples can change the number of
clusters, and so the entropy, without changing a single entailment judgment.

An earlier version of this work grouped answers by the connected components of
the mutual-entailment graph, using union-find. That is the natural reading of
``group answers that mean the same thing'', it is order-independent, and it
merges transitively. It is also not what the code does, and the two procedures
give different cluster counts on the same entailment matrix whenever the
relation is not transitive, which for a noisy neural model is common.
The point of this thesis is to measure the method as published, so the anchor
scan is used throughout and its order dependence is carried as a property of
the method under study. Algorithm~\ref{alg:clustering} states it
exactly.

\begin{algorithm}
\caption{Bidirectional-entailment clustering by order-dependent anchor scan,
  reproducing \texttt{get\_semantic\_ids} in the released
  code~\cite{kossen_semantic_uncertainty_code_2024}.}
\label{alg:clustering}
\KwIn{answers $y_1,\ldots,y_M$ in generation order; entailment model
  $\mathcal{M}(\cdot,\cdot)$ returning entailment, neutral, or
  contradiction (queries are prefixed with the question in the
  question-conditioned configuration)}
\KwOut{one cluster identifier per answer}
$c_i \gets \textsc{unassigned}$ for $i=1,\ldots,M$\;
$k \gets 0$\tcp*{next cluster identifier}
\For{$i \gets 1$ \KwTo $M$}{
  \If{$c_i$ is \textsc{unassigned}}{
    $c_i \gets k$\tcp*{$y_i$ becomes an anchor}
    \For{$j \gets i+1$ \KwTo $M$}{
      $a \gets \mathcal{M}(y_i, y_j)$\tcp*{does $y_i$ entail $y_j$?}
      $b \gets \mathcal{M}(y_j, y_i)$\tcp*{and $y_j$ entail $y_i$?}
      \If{$a$ is entailment \textbf{and} $b$ is entailment}{
        $c_j \gets k$\tcp*{assigned even if $y_j$ already had a cluster}
      }
    }
    $k \gets k + 1$\;
  }
}
\Return $(c_1,\ldots,c_M)$\;
\end{algorithm}

\section{Discrete semantic entropy}
\label{sec:discrete-se}

After clustering, let $n_k$ be the number of answers in cluster $c_k$. Its
empirical probability is $\hat{p}_k = n_k/M$, and the discrete semantic entropy
is
\[
  \widehat{H}_{\mathrm{sem}}(x)
  \;=\; -\sum_{k=1}^{K} \hat{p}_k \log \hat{p}_k .
\]
A low score means the samples mostly agree on one meaning; a high score means
they spread across several. The maximum is $\log K \le \log M$. The experiments
use $M=10$, so at most $\ln 10 \approx 2.30$; the worked example below uses
$M=5$, so at most $\ln 5 \approx 1.61$.

\paragraph{The estimator family.}\label{sec:estimator-family}
The grid reports this count-based estimator, the probability-weighted variant
of Farquhar et al.~\cite{semantic_entropy_nature_2024,%
kossen_semantic_uncertainty_code_2024}, and two clustering-free scores. Two
axes separate them: whether answers are grouped by meaning, and whether the
score reads sequence probabilities or only counts.

\begin{center}
\small
\resizebox{\linewidth}{!}{%
\begin{tabular}{@{}lll@{}}
\toprule
 & Count-based & Probability-weighted \\
\midrule
No meaning judgment & \makecell[l]{naive sample entropy (raw strings),\\
  surface entropy (normalized strings)} &
  \makecell[l]{naive entropy\\(the paper's naive baseline)} \\
Meaning judgment & \makecell[l]{discrete semantic entropy\\
  (the paper's discrete approximation)} &
  \makecell[l]{full semantic entropy\\(the paper's primary estimator)} \\
\bottomrule
\end{tabular}}
\end{center}

Two of the four ignore meaning and read only the strings: \emph{naive sample
entropy} counts the raw strings, \emph{surface entropy} counts them after
normalization. A third, the paper's \emph{naive entropy}, ignores meaning
differently: it does no grouping at all and is the mean negative
length-normalized sequence log-likelihood over the $M$ samples, so it reads how
confident the model was rather than how often an answer recurred. The two
semantic scores group first: \emph{discrete semantic entropy} counts the
groups, \emph{full semantic entropy} weights them by the summed likelihood of
their members. When the samples are equally likely the full version reduces
exactly to the discrete one, which is a useful check on the implementation.

The distinction between naive entropy and the two string scores is what RQ2
rests on, and it is easy to lose. Naive entropy differs from discrete semantic
entropy in \emph{two} ways: it reads likelihoods rather than counts, and it does
no grouping. Surface entropy differs in \emph{one}: it counts exactly as the
discrete estimator does, but groups by string identity rather than by meaning.
So a comparison against naive entropy cannot separate the contribution of the
grouping from the contribution of counting. A comparison against surface
entropy can.

Algorithm~\ref{alg:scoring} gives the procedure. Note the ordering: clustering
consumes the raw answers, and normalization is applied only on the branch that
produces the comparator.

\begin{algorithm}
\caption{Semantic-entropy scoring from sampled answers, with the
  clustering-free comparator computed from the same samples.}
\label{alg:scoring}
\KwIn{raw sampled answers $y_1,\ldots,y_M$ for prompt $x$; optional
  length-normalized sequence log-likelihoods $\ell_1,\ldots,\ell_M$;
  entailment model $\mathcal{M}$}
\KwOut{discrete semantic entropy $\widehat{H}_{\mathrm{sem}}$, surface
  entropy $\widehat{H}_{\mathrm{lex}}$, and, when log-likelihoods are
  available, full semantic entropy $\widehat{H}^{\mathrm{p}}_{\mathrm{sem}}$}
$(c_1,\ldots,c_M) \gets \textsc{cluster}(y_1,\ldots,y_M;\mathcal{M})$\tcp*{Algorithm~\ref{alg:clustering}, on raw answers}
let $K$ be the number of distinct clusters\;
\For{$k \gets 1$ \KwTo $K$}{
  $n_k \gets |\{\,i : c_i = k\,\}|$\;
  $\hat{p}_k \gets n_k / M$\;
}
$\widehat{H}_{\mathrm{sem}} \gets -\sum_{k=1}^{K}\hat{p}_k \log \hat{p}_k$\tcp*{discrete semantic entropy}
\tcp{comparator: normalize, then count distinct strings}
$u_i \gets \textsc{normalize}(y_i)$ for $i=1,\ldots,M$\;
$\hat{q}_s \gets |\{\,i : u_i = s\,\}| / M$ for each distinct string $s$\;
$\widehat{H}_{\mathrm{lex}} \gets -\sum_{s}\hat{q}_s \log \hat{q}_s$\tcp*{surface entropy}
\If{sequence log-likelihoods are available}{
  $\log P_k \gets \operatorname{logsumexp}_{\,i:\,c_i=k}\ \ell_i$ for each $k$\;
  $Z \gets \operatorname{logsumexp}_{\,k}\ \log P_k$\tcp*{renormalize cluster masses}
  $q_k \gets \exp(\log P_k - Z)$ for each $k$\;
  $\widehat{H}^{\mathrm{p}}_{\mathrm{sem}} \gets -\sum_{k=1}^{K} q_k \log q_k$\tcp*{full semantic entropy}
}
\Return $\widehat{H}_{\mathrm{sem}}$, $\widehat{H}_{\mathrm{lex}}$ (and
  $\widehat{H}^{\mathrm{p}}_{\mathrm{sem}}$ if computed)\;
\end{algorithm}

\section{Worked example}
\label{sec:worked_example}

The following example makes the pipeline concrete. Consider the factual
question $x = $ ``What is the capital of Australia?''. The language model is
sampled $M=5$ times under one decoding configuration, drawing a fresh
completion each time. That yields five different but plausible answers:

\begin{center}
\begin{tabular}{cl}
\hline
Index & Sampled answer $y_i$ \\
\hline
1 & \emph{Canberra} \\
2 & \emph{Canberra is the capital city of Australia} \\
3 & \emph{Sydney} \\
4 & \emph{Melbourne} \\
5 & \emph{Canberra} \\
\hline
\end{tabular}
\end{center}

The example is intentionally small and intentionally a case in which the model
is mostly right but somewhat uncertain, which makes it a useful check on the
pipeline. The decoding settings of the recorded experiments are given in
Section~\ref{sec:sampling-settings}.

\paragraph{Step 1: Bidirectional entailment.}
The entailment model is applied to each ordered pair $(y_i, y_j)$.
Table~\ref{tab:entailment_example} shows the resulting binary decisions; a
cell $(i,j)$ equals 1 if the row answer entails the column answer.

\begin{table}[ht]
\thisfloatsetup{capposition=above}
\centering
\caption{Pairwise entailment matrix for the worked example.
         A cell $(i,j)=1$ means $y_i$ entails $y_j$.}
\label{tab:entailment_example}
\begin{tabular}{c|ccccc}
       & $y_1$ & $y_2$ & $y_3$ & $y_4$ & $y_5$ \\\hline
$y_1$ &  --   &   1   &   0   &   0   &   1   \\
$y_2$ &   1   &  --   &   0   &   0   &   1   \\
$y_3$ &   0   &   0   &  --   &   0   &   0   \\
$y_4$ &   0   &   0   &   0   &  --   &   0   \\
$y_5$ &   1   &   1   &   0   &   0   &  --   \\
\end{tabular}
\end{table}

\paragraph{Step 2: Semantic clustering.}
Two answers belong to the same cluster when both off-diagonal cells are 1.
Reading from Table~\ref{tab:entailment_example}: $y_1 \sim y_2$,
$y_1 \sim y_5$ and $y_2 \sim y_5$, while $y_3$ and $y_4$ each stand alone.

It is worth tracing the anchor scan of Algorithm~\ref{alg:clustering} rather
than assuming the answer, since the scan and a connected-components grouping
do not agree in general. Answer $y_1$ is unassigned, so it anchors cluster~0;
the inner loop finds $y_2$ and $y_5$ mutually entailed with it and assigns
both to cluster~0, and rejects $y_3$ and $y_4$. The scan moves on: $y_2$ is
already assigned, $y_3$ is not and anchors cluster~1, matching nothing;
$y_4$ anchors cluster~2, matching nothing; $y_5$ is already assigned. The
entailment relation here happens to be transitive on
$\{y_1,y_2,y_5\}$, so no reassignment occurs and the scan returns the same
partition a connected-components grouping would. This yields $K=3$ clusters:
\[
  c_1 = \{y_1, y_2, y_5\},\quad
  c_2 = \{y_3\},\quad
  c_3 = \{y_4\}.
\]
Section~\ref{sec:order_dependence_example} shows a case in which that
agreement fails.

\paragraph{Step 3: Empirical probabilities.}
\[
  \hat{p}_1 = \tfrac{3}{5} = 0.60,\quad
  \hat{p}_2 = \tfrac{1}{5} = 0.20,\quad
  \hat{p}_3 = \tfrac{1}{5} = 0.20.
\]

\paragraph{Step 4: Discrete semantic entropy.}
\begin{align*}
  \widehat{H}_{\mathrm{sem}}(x)
  &= -\bigl(0.60\ln 0.60 + 0.20\ln 0.20 + 0.20\ln 0.20\bigr) \\
  &= -(-0.306 - 0.322 - 0.322) \\
  &\approx 0.95 \text{ nats.}
\end{align*}

\paragraph{The clustering-free comparator.}
This is not a step of the semantic-entropy method but the comparator RQ2
turns on, computed from the same five samples. It treats each distinct
normalized string as its own outcome. With strings
$\{$\emph{Canberra}$\times 2$, \emph{Canberra is the capital city of
Australia}$\times 1$, \emph{Sydney}$\times 1$, \emph{Melbourne}$\times 1\}$,
the surface entropy is
\begin{align*}
  \widehat{H}_{\mathrm{lex}}(x)
  &= -\bigl(0.40\ln 0.40 + 0.20\ln 0.20 + 0.20\ln 0.20 + 0.20\ln 0.20\bigr) \\
  &\approx 1.33 \text{ nats.}
\end{align*}

In this example the count over normalized strings and the count over raw
strings coincide, because no two of the five samples differ only in case or
punctuation. Chapter~\ref{ch:results} keeps them apart as \emph{surface
entropy} and \emph{naive sample entropy}, since on the recorded data they do
not always agree.

The surface score is higher (1.33 against 0.95) because it treats the
paraphrase ``Canberra is the capital city of Australia'' as a different answer
from ``Canberra'', inflating the apparent uncertainty. After semantic
clustering the correct answer class accounts for 60\,\% of the samples, and
the entropy is correspondingly lower. This is the mechanism the method is
named for, and the example is constructed to display it. Whether it produces a
measurable advantage on real data, once both scores are computed over tens of
thousands of records, is the question Chapter~\ref{ch:results} settles;
the example establishes what is being claimed, not that the claim holds.

\subsection{When the grouping depends on the order}
\label{sec:order_dependence_example}

The example above is well behaved: its entailment relation happens to be
transitive, so the anchor scan of Algorithm~\ref{alg:clustering} and a
connected-components grouping return the same three clusters. That is not
guaranteed, and when it fails the consequences are larger than they might
appear. The following case shows why.

Take four sampled answers to the same question:

\begin{center}
\begin{tabular}{cl}
\hline
Index & Sampled answer \\
\hline
$a$ & \emph{Canberra} \\
$b$ & \emph{Canberra, the capital} \\
$c$ & \emph{The capital city} \\
$d$ & \emph{Sydney} \\
\hline
\end{tabular}
\end{center}

Suppose the entailment model judges $a \leftrightarrow b$ mutually entailed
and $b \leftrightarrow c$ mutually entailed, but \emph{not}
$a \leftrightarrow c$: ``Canberra'' on its own does not entail ``the capital
city'', and vice versa, without the question supplying the link that $b$
states explicitly. This is not a contrived failure. It is the ordinary
behavior of a relation that is symmetric by construction but not transitive,
and a neural entailment model has no mechanism enforcing transitivity.

Grouping by connected components merges all three through the bridge $b$,
giving $\{a,b,c\},\{d\}$ regardless of the order the samples arrived in. The
anchor scan does not, and what it returns depends on that order:

\begin{center}
\begin{tabular}{lll}
\hline
Order presented & Resulting clusters & $\widehat{H}_{\mathrm{sem}}$ \\
\hline
$a,b,c,d$ & $\{a,b\},\{c\},\{d\}$      & $1.04$ nats \\
$b,a,c,d$ & $\{a,b,c\},\{d\}$          & $0.56$ nats \\
$a,c,b,d$ & $\{a\},\{b,c\},\{d\}$      & $1.04$ nats \\
\hline
\end{tabular}
\end{center}

Three permutations of the same four answers give the same entailment judgments
throughout, but two different partitions. One of them assigns nearly twice the
entropy of the other.

The third row is the more striking case, and it is the reassignment rule at
work. Answer $a$ anchors first and claims $b$. The scan then reaches $c$,
which is still unassigned, so $c$ becomes an anchor in its own right and
claims $b$ a second time and overwrites $b$'s existing identifier. Answer $a$ is
left in a cluster by itself, even though it is mutually entailed with an answer
it originally grouped with.

Two points follow. First, the reported entropy is a property of the generation
order as well as of the entailment judgments, which is worth knowing when the
estimator is described as measuring ``the number of distinct meanings''.
Second, the effect is invisible in aggregate: the partition is well defined
and the entropy is finite, so nothing downstream signals that a different
ordering would have produced a different number.

Whether this matters empirically is a separate question from whether it
occurs, and it is settled by measurement rather than argument. The
disagreement between entailment backends reported in
Section~\ref{sec:clustering-sensitivity} bounds how much the clustering step
moves the results, and order effects are one component of that. This example
establishes only that the estimator as published is order-dependent. A reader
who assumes otherwise, as the natural reading of ``cluster by meaning''
invites, would be mistaken.

\section{Method contract}
\label{sec:method-contract}

The algorithm itself needs little. It takes a prompt $x$ with $M$ sampled
answers drawn under a fixed decoding configuration, and optionally their
length-normalized sequence log-likelihoods. It groups the \emph{raw} answers
with the configured equivalence relation, each entailment query prefixed with
$x$ in the question-conditioned configuration, applying no normalization first.
It returns the cluster assignments and the discrete semantic entropy over
cluster frequencies, plus the full semantic entropy when log-likelihoods are
supplied. From the same samples it also applies the dataset's answer-cleaning
and returns the surface entropy over distinct normalized strings.

Correctness labelling and evaluation consume those outputs but are not part of
the method. That separation is what lets
Chapter~\ref{ch:experimental_design} apply four entailment backends and three
correctness rules to one stored set of generations.

\section{Related estimators evaluated alongside}
\label{sec:related-estimators}

The remaining estimators are published baselines, defined here only as far as
Chapter~\ref{ch:results} needs. One of them, the Semantic Entropy Probe,
addresses the main method's cost: discrete semantic entropy needs $M$ samples
per question and up to $M^2$ pairwise entailment checks. The other two, P(True)
and embedding regression, are simple and widely used.

All three read a single answer per question rather than the ten the entropy
estimators use. That answer is called the \emph{most-likely answer} throughout,
and it is also what correctness is scored on. It is not a greedy decode:
following the released code, it is a sample drawn at temperature $0.1$, under a
token budget four times larger than the samples get.
Section~\ref{sec:sampling-settings} gives the details.

\subsection{Semantic Entropy Probes}
\label{sec:sep-method}

A Semantic Entropy Probe (SEP)~\cite{kossen_semantic_entropy_probes_2024}
attacks the cost problem from the other end: it reads uncertainty off the
model's internal state, so at test time one forward pass is enough, with no
multi-sampling and no entailment comparisons.

For each training question, the most-likely answer supplies a hidden state,
read at the last generated token as fixed by the stop sequence. Separately, ten
sampled responses are grouped by meaning to compute discrete semantic entropy.
That entropy is turned into a high/low label, and a logistic-regression
classifier is trained to predict it from the hidden state. Correctness labels
are not used; the accuracy probe below is the variant trained on those. At test
time, generate one answer, read the same hidden state, apply the classifier,
and use its output as the uncertainty score.

The probe is trained separately per dataset, model and seed on its own prompts
and scored on a disjoint set, so its numbers are a genuine held-out test. The
feature choices are fixed in advance rather than tuned on the evaluation set.
One detail follows the code rather than the natural reading: the hidden state is
taken at position $n_{\mathrm{generated}}-1$, where $n_{\mathrm{generated}}$
comes from where the stop sequence cut the answer, and no end-of-sequence token
has to be present. The probe is far cheaper than full semantic entropy, but
limited to what is linearly readable from one hidden state.

\subsection{P(True) and embedding regression}
\label{sec:other-estimators}

\textbf{P(True)} asks the model to grade itself. It is shown the question, a
set of brainstormed alternative answers, and its own most-likely answer, and is
asked whether that answer is true, with the options presented as letters; a
twenty-shot prefix from the training split comes first. What the released code
records is the raw log-probability the model puts on the token for the ``true''
option, not a probability renormalized over the two options, which is what the
published description suggests. This thesis follows the code, and
Section~\ref{sec:deviation-register} records the difference. The score is
oriented like every other score here, larger meaning riskier. P(True) needs no
extra samples and no probe training, which makes it a natural cheap reference.

\textbf{Embedding regression} is a close cousin of the SEP probe: the same kind
of linear model on the same hidden states of the same most-likely answers, but
trained on \emph{judged correctness} instead of a semantic-entropy threshold.
It is reported in two forms. \emph{In-distribution} means trained and tested on
the same task. \emph{Out-of-distribution} means trained on the pooled prompts of
the other two short-answer tasks and scored on the held-out one. Together they
measure how well such a probe transfers between datasets.
Chapter~\ref{ch:results} calls it the \emph{accuracy probe}, since its training
target is what distinguishes it from SEP.

\chapter{Experimental Design}

TODO: Define datasets, models, sampling settings, clustering, correctness labels, and evaluation metrics.

\chapter{Results}
\label{ch:results}

This chapter reports what the experiments found. It is written to be
readable without prior knowledge of the method: every table and figure
is introduced in plain words before any numbers are discussed, and a set
of worked, by-hand examples (Section~\ref{sec:hand-audits}) shows what
the detector actually does on individual questions before the
aggregate scores are presented.

\section{How to read the numbers}
\label{sec:how-to-read}

The task in every experiment is the same. The language model answers a
question; sometimes it is right and sometimes it is wrong. A
\emph{detector} looks at the answer (or at several sampled answers) and
outputs a single number that is meant to be high when the answer is
unreliable. A good detector gives higher numbers to the wrong answers
than to the right ones, so that a user could decline the riskiest
answers and keep the safe ones.

Two summary measures are used throughout.

\begin{itemize}
  \item \textbf{AUROC} (area under the ROC curve). Pick one wrong answer
    and one right answer at random; the AUROC is the probability that
    the detector scored the wrong one higher. \textbf{0.5 means the
    detector is useless} (no better than a coin flip), and \textbf{1.0
    means it is perfect}. Values around 0.7--0.9 are useful in practice.
    Because AUROC only depends on the \emph{ranking} of scores, the unit
    of the entropy (nats vs.\ bits) does not affect it.
  \item \textbf{AURAC} (area under the rejection--accuracy curve). If we
    throw away the answers the detector is most worried about, the
    accuracy of what remains should go up. AURAC summarizes how quickly
    accuracy improves as we discard flagged answers; higher is better.
    Unlike AUROC, it also depends on how accurate the model was to begin
    with, so AURAC values are only comparable within the same task and
    model.
\end{itemize}

Throughout, the positive class is the \emph{incorrect} answer, so
``higher score $=$ more likely wrong''. All headline numbers are
averaged over three independent runs (random seeds 42, 43, 44); where a
spread is shown it is the standard deviation across those three runs.
The detectors compared are the ones defined in
Chapter~\ref{ch:semantic_entropy}: the cheap surface baselines
(counting distinct answer strings), discrete semantic entropy with
meaning-based clustering, the probability-weighted variants, Kernel
Language Entropy, the hidden-state probe (SEP), \mbox{P(True)}, and
embedding regression.

\section{The headline picture}
\label{sec:headline}

Figure~\ref{fig:regime-split} is the single most important result of the
thesis, and the rest of the chapter unpacks it. It compares two
detectors across the four tasks and three models: the \emph{cheap}
surface score, which just counts how many of the sampled answers are
different strings, and the \emph{meaning-aware} semantic entropy, which
first groups answers that mean the same thing and then measures the
spread over meanings.

\begin{figure}[ht]
\centering
\includegraphics[width=\textwidth]{regime_split.png}
\caption{Surface-form versus meaning-aware detection across tasks
  (AUROC, higher is better; the dashed line is the useless ``coin-flip''
  level of 0.5). On the three short-answer tasks the two scores are
  almost the same height: the cheap string-counting baseline is already
  as good as meaning-aware clustering. On the long-form biography task
  the surface bar drops to the chance line --- it stops working
  entirely --- while the meaning-aware bar stays well above it. This
  short-versus-long contrast is the central finding.}
\label{fig:regime-split}
\end{figure}

The pattern is the same for all three models. On short answers the two
bars are nearly equal, so the expensive meaning-based machinery buys
almost nothing over simply noticing that the model wrote different
strings. On long-form biographies the cheap baseline collapses to the
chance line while the meaning-aware score survives. \emph{Where you need
semantic clustering is precisely the long-form regime; on short answers
it is close to redundant.} The remainder of the chapter establishes this
quantitatively and explains why it happens.

\section{Worked examples: what the detector sees}
\label{sec:hand-audits}

Before the aggregate tables, it helps to look at three real records from
the recorded runs (canonical configuration, temperature $1.0$). Each one
shows the ten sampled answers, how the meaning-based clustering grouped
them, and the resulting scores. These by-hand audits make concrete what
the averages later summarize.

\subsection{Example 1 --- meaning-based grouping lowers a falsely high
score}
\label{sec:audit-merge}

\textbf{Question (TriviaQA, Llama-3.1-8B):} ``According to Rudyard
Kipling, what were the two impostors to meet and treat the same day?''
\textbf{Gold answer:} \emph{Triumph and Disaster}. The ten sampled
answers were:

\begin{quote}\small
\emph{Success and Failure.} $\;\cdot\;$ \emph{Triumph and Disaster}
$\;\cdot\;$ \emph{Triumph and defeat.} $\;\cdot\;$ \emph{Triumph and
disaster.} $\;\cdot\;$ \emph{Triumph and defeat.} $\;\cdot\;$
\emph{Success and Failure} $\;\cdot\;$ \emph{Success and Failure.}
$\;\cdot\;$ \emph{Triumph and Disaster} $\;\cdot\;$ \emph{Tyrant and a
dreamer.} $\;\cdot\;$ \emph{Success and Failure.}
\end{quote}

The surface baseline sees many \emph{different strings} (``Triumph and
Disaster'', ``Triumph and defeat.'', ``Success and Failure'', and so on)
and therefore reports high uncertainty: surface entropy $1.28$ nats. The
meaning-based step asks, for each pair, whether the two phrases say the
same thing in context; it judges the various ``triumph/defeat'' and
``success/failure'' phrasings to be mutually entailing and collapses
nine of the ten samples into a single meaning, leaving only the odd
``Tyrant and a dreamer'' on its own. The semantic entropy is then only
$0.33$ nats --- a much more honest reflection of the fact that the model
basically settled on one idea. The answer was graded correct. This is
the mechanism by which clustering is \emph{supposed} to help: it removes
uncertainty that is only about wording.

\subsection{Example 2 --- a confident wrong answer the method cannot
catch}
\label{sec:audit-artifact}

\textbf{Question (TriviaQA, Llama-3.1-8B):} ``Who was the director of the
CIA from 1976--81?'' \textbf{Gold answer:} \emph{George Bush}. All ten
samples were the same: \emph{Stansfield Turner}. Because every sample
agrees, the spread over meanings is zero (semantic entropy $0.00$), so
the detector reports maximum confidence --- yet the answer is graded
incorrect.

This single record illustrates two important limits at once. First, it
is a genuine \emph{confident systematic error}: the model is wrong in the
same way every time, and no consistency-based detector can flag it,
because there is no inconsistency to measure. Second, it is partly a
\emph{labelling artefact}: Stansfield Turner did head the CIA for most of
1977--1981, so the ``wrong'' label here reflects an ambiguous gold answer
as much as a model error. Both effects --- the blind spot for confident
errors and the noise in automatic correctness labels --- recur in the
aggregate numbers and are returned to in the discussion
(Chapter~\ref{ch:discussion}).

\subsection{Example 3 --- why the surface baseline dies on biographies}
\label{sec:audit-bio}

\textbf{Question (Biography, Mistral-7B):} ``Tell me a bio of Edmond
Pourchot.'' The model produced ten multi-sentence paragraphs. One began:
\emph{``Edmond Pourchot, born on March 2, 1896, in Gap, France, was a
prominent figure in the field of French aviation\ldots''} (Edmond
Pourchot was in fact a seventeenth-century French philosopher; the
paragraph is fabricated.) All ten paragraphs were different strings, so
the surface baseline reports the \emph{maximum possible} value every
time: $\ln 10 = 2.30$ nats. Across the whole biography set the average
surface entropy is $2.303$ --- pinned at the ceiling --- which is exactly
why it carries no information and lands at the chance line in
Figure~\ref{fig:regime-split}: a score that is always maximal cannot rank
anything.

Meaning-based clustering escapes this trap. It groups the paragraphs by
what they claim rather than by their exact wording, so its score varies
from record to record (average semantic entropy $1.73$, not pinned at the
ceiling). For this fabricated biography the ten paragraphs genuinely
disagree, so the semantic entropy is high and the answer is correctly
flagged as risky. This is the long-form mechanism behind the headline
figure.

\section{Short-answer results}
\label{sec:short-answer-results}

We begin with the three short-answer tasks (TriviaQA, NQ-Open, SVAMP).
Table~\ref{tab:auroc-overview} gives the bird's-eye view: discrete
semantic entropy, the full probability-weighted estimator, and the
hidden-state probe, for all three models and four tasks.

\begin{table}[ht]
\thisfloatsetup{capposition=above}
\centering
\caption{Selective-prediction AUROC overview (seed-mean over 42/43/44).
  Short-answer tasks on the left, biography on the right. On short
  answers all detectors land in a similar, useful band; on biography the
  picture changes sharply (Section~\ref{sec:bio-results}).}
\label{tab:auroc-overview}
\input{includes/generated_overview.tex}
\end{table}

\begin{figure}[ht]
\centering
\includegraphics[width=\textwidth]{multiseed_auroc.png}
\caption{Per-task AUROC for the main detectors, averaged over three
  seeds (error bars are the seed spread). The detectors cluster
  together on every short-answer task: no single method dominates, and
  the cheap surface baseline sits inside the same band as the
  meaning-aware scores.}
\label{fig:multiseed-auroc}
\end{figure}

The detailed per-task tables follow.
Tables~\ref{tab:auroc-triviaqa}, \ref{tab:auroc-nqopen}
and~\ref{tab:auroc-svamp} list every detector for TriviaQA, NQ-Open and
SVAMP respectively.

\begin{table}[ht]
\thisfloatsetup{capposition=above}
\centering
\caption{TriviaQA AUROC by detector and model (mean $\pm$ std over
  seeds). Canonical DeBERTa-v3-large clustering for the entropy methods.}
\label{tab:auroc-triviaqa}
\input{includes/generated_auroc_triviaqa.tex}
\end{table}

\begin{table}[ht]
\thisfloatsetup{capposition=above}
\centering
\caption{NQ-Open AUROC by detector and model (mean $\pm$ std over seeds).}
\label{tab:auroc-nqopen}
\input{includes/generated_auroc_nqopen.tex}
\end{table}

\begin{table}[ht]
\thisfloatsetup{capposition=above}
\centering
\caption{SVAMP AUROC by detector and model (mean $\pm$ std over seeds).}
\label{tab:auroc-svamp}
\input{includes/generated_auroc_svamp.tex}
\end{table}

Three things stand out on short answers. First, all of the
sampling-based detectors work: AUROC sits in roughly the
$0.74$--$0.88$ band, which is consistent with the original semantic
entropy paper under the matched protocol. Second, the
\emph{surface/exact-match baseline is the strongest single score in most
cells} --- for example on TriviaQA it reaches $0.819$ (Mistral),
$0.838$ (Llama-3.1-8B) and $0.802$ (Llama-3.1-70B), slightly ahead of
the meaning-aware $0.799$, $0.814$ and $0.794$. It is the single
strongest score in six of the nine short-answer cells (all of TriviaQA
and NQ-Open); only on SVAMP does the probability-weighted full estimator
edge ahead. Third, discrete semantic
entropy and the full probability-weighted estimator track each other
closely, confirming that on these tasks the cheap count-based version
loses almost nothing relative to the version that uses the model's own
probabilities. The reason, made concrete in
Example~\ref{sec:audit-merge}, is that on short answers most of the
useful grouping is already done by \emph{normalizing} the strings;
meaning-based entailment then has little left to add.

\section{Does meaning-based clustering beat the surface baseline?}
\label{sec:significance}

Because the surface and semantic scores are so close on short answers,
the honest question is whether the difference is real or just noise.
Table~\ref{tab:significance-vs-surface} reports paired bootstrap
comparisons: for each task and model it asks whether a given detector
beats the surface baseline by more than sampling noise would explain.

\begin{table}[ht]
\thisfloatsetup{capposition=above}
\centering
\caption{Paired-bootstrap comparison against the surface baseline.
  A confidence interval that contains zero means ``not significantly
  different from the cheap baseline.''}
\label{tab:significance-vs-surface}
\input{includes/generated_significance.tex}
\end{table}

On the short-answer tasks, \emph{no method beats the surface baseline by
a statistically significant margin}: wherever a method's margin over
surface entropy is large enough to be significant, it is the surface
baseline that comes out ahead, never the other way round. In other
words, once you have normalized the answer strings, paying for an NLI
model to cluster them does not reliably improve short-answer
hallucination detection.
This is a negative result, and an important one: it says the expensive
part of the method earns its keep only in the long-form setting, which
the next sections examine.

\section{Clustering backend sensitivity}
\label{sec:clustering-sensitivity}

The meaning-based step needs a model to decide whether two answers are
equivalent. Does the choice of that model matter?
Table~\ref{tab:clustering-sensitivity} compares five backends ---
exact-match, two sizes of DeBERTa NLI, and a large Qwen judge ---
holding everything else fixed.

\begin{table}[ht]
\thisfloatsetup{capposition=above}
\centering
\caption{Discrete semantic entropy AUROC by clustering backend
  (mean $\pm$ std over seeds). ``--'' marks a backend that was not run
  for that task (the larger NLI judges were applied to short answers
  only as a sensitivity check).}
\label{tab:clustering-sensitivity}
\input{includes/generated_clustering_sensitivity.tex}
\end{table}

\begin{figure}[ht]
\centering
\includegraphics[width=0.85\textwidth]{bio_clustering_sensitivity.png}
\caption{Biography AUROC for discrete semantic entropy under different
  clustering backends. On long-form inputs the backend is
  result-determining: the same samples yield anything from chance-level
  to clearly useful detection depending only on how equivalence is
  judged.}
\label{fig:bio-clustering-sensitivity}
\end{figure}

On short answers the backend barely matters --- exact-match, base NLI,
large NLI and the Qwen judge land within a few points of each other,
which is unsurprising given that normalization already does the work. On
biographies the choice is decisive: the Qwen judge clusters meanings more
reliably than the smaller DeBERTa models and lifts AUROC accordingly.
This is the long-form counterpart to the short-answer negative result:
when clustering matters at all, \emph{which} clusterer you use matters a
great deal.

\section{Long-form biography results}
\label{sec:bio-results}

The biography task is the controlled long-form case: the same
sample-and-cluster pipeline, but the model now writes a paragraph rather
than a phrase. Table~\ref{tab:auroc-bio} lists every detector for all
three generators.

\begin{table}[ht]
\thisfloatsetup{capposition=above}
\centering
\caption{Biography (whole-paragraph) AUROC by detector and model
  (mean $\pm$ std over seeds), canonical DeBERTa-v3-large clustering.
  The two count-based surface scores (naive \emph{sample} entropy and
  surface entropy) collapse to chance ($\approx 0.50$) because long
  paragraphs almost never repeat as identical strings; the
  probability-weighted naive entropy, the meaning-aware methods, and the
  hidden-state probe all recover a real signal.}
\label{tab:auroc-bio}
\input{includes/generated_auroc_bio.tex}
\end{table}

The contrast with short answers is stark. The two count-based surface
scores --- naive \emph{sample} entropy and surface entropy --- sit at
$\approx 0.50$, pure chance, exactly as Example~\ref{sec:audit-bio}
predicted: every long paragraph is a distinct string, so these scores
are permanently pinned at their maximum and can rank nothing. Every
detector that escapes raw string-counting recovers a real signal.
Meaning-aware discrete semantic entropy reaches $0.708$ (Mistral-7B),
$0.618$ (Llama-3.1-8B) and $0.612$ (Llama-3.1-70B), and the Qwen-judge
clustering variant is stronger still ($0.736 / 0.685 / 0.666$),
consistent with the backend sensitivity above. So the value of moving
from surface comparison to meaning comparison is not marginal here --- it
is the difference between a useless detector and a working one.

Two detectors that need no clustering at all are in fact the strongest on
biography. Probability-weighted naive entropy --- which keeps each
distinct paragraph but weights it by the model's own sequence
probability --- is the single best score on the two smaller models
($0.781$ Mistral, $0.769$ Llama-3.1-8B), and the self-evaluation
baseline \mbox{P(True)} is the best on the largest ($0.777$). The lesson
is that on long text the model's \emph{own probability signal} carries
much of the confabulation information, and the hard-clustering step only
partially recovers it.

One regularity is worth flagging. The semantic-entropy AUROC
\emph{decreases} as the generator gets larger ($0.708 \to 0.618 \to
0.612$ from 7B to 70B). The discussion returns to this, but the short
explanation is that larger models confabulate more \emph{consistently}:
when a big model invents a fact, its ten samples tend to agree on the
invention, leaving little spread for a consistency-based detector to
find.

\section{Supervised baselines and hidden-state probes}
\label{sec:probes-results}

Semantic entropy needs several samples per question. Two cheaper-at-test
alternatives read the model's internal state from a single pass: the
hidden-state probe (SEP) and embedding regression.
Table~\ref{tab:embedding-regression} reports them alongside
\mbox{P(True)}.

\begin{table}[ht]
\thisfloatsetup{capposition=above}
\centering
\caption{Probe and supervised-baseline AUROC (mean $\pm$ std over
  seeds): hidden-state SEP probe, embedding regression
  (in-distribution and out-of-distribution), and \mbox{P(True)}.}
\label{tab:embedding-regression}
\input{includes/generated_embedding_regression.tex}
\end{table}

The probes are competitive, and on biography they are strong: the SEP
probe reaches $0.728$ (8B) and $0.731$ (70B), \emph{above} every
meaning-based clustering method on those models, and \mbox{P(True)} is
the single best biography detector on the 70B model at $0.777$. (Only
the probability-weighted naive entropy, noted above, beats the probe on
the two smaller models.) This matters for deployment, because a
single-pass probe is far cheaper than drawing ten samples and running
pairwise entailment. The catch, shown next, is that the probe's
advantage depends on the decoding temperature.

\section{Sampling-temperature sensitivity}
\label{sec:temperature-sensitivity}

Every result so far used the model's default sampling temperature. Since
sampling-based detectors depend on getting \emph{diverse} samples, the
decoding temperature ought to matter --- and it does, in a way that has a
clean practical lesson. Table~\ref{tab:temp-sweep} and
Figure~\ref{fig:temp-sweep} sweep the temperature from $0.1$ (nearly
deterministic) to $1.0$.

\begin{table}[ht]
\thisfloatsetup{capposition=above}
\centering
\caption{AUROC versus sampling temperature (seed-mean), Llama-3.1-70B.
  Sampling-based entropy rises with temperature; the single-pass probe
  falls.}
\label{tab:temp-sweep}
\input{includes/generated_temp_sweep_auroc.tex}
\end{table}

\begin{figure}[ht]
\centering
\includegraphics[width=\textwidth]{auroc_vs_temperature.png}
\caption{AUROC as a function of sampling temperature, per task and
  model. The sampling-based scores (semantic entropy, surface entropy)
  climb as temperature increases, because more diverse samples reveal
  more disagreement to measure. The hidden-state probe moves the
  opposite way. The two families therefore \emph{cross over} between
  temperature $0.7$ and $1.0$ on short answers; on biography the curves
  are flat.}
\label{fig:temp-sweep}
\end{figure}

At low temperature the model repeats nearly the same answer every time,
so there is little diversity for a sampling-based method to work with and
semantic entropy is weak (TriviaQA/70B: AUROC $0.54$ at $T{=}0.1$). As
temperature rises the samples spread out and the same method becomes
strong ($0.78$ at $T{=}1.0$). The hidden-state probe does the reverse ---
it is best at low temperature and degrades as decoding gets noisier ---
so the two families \emph{cross over}. The practical lesson is that there
is no single best detector: at low decoding temperature, use the probe;
at high temperature, use sampling-based entropy. On biography the curves
are flat, so this trade-off is specific to the short-answer regime.

\section{Claim-level long-form detection: FactualBio}
\label{sec:factualbio-results}

The whole-paragraph biography score asks ``is this paragraph reliable?''
A finer question is ``which individual claims in the paragraph are
reliable?'' The FactualBio extension breaks each biography into
atomic claims and applies the same meaning-based idea at the claim level.
Table~\ref{tab:factualbio} reports the result.

\begin{table}[ht]
\thisfloatsetup{capposition=above}
\centering
\caption{Claim-level FactualBio AUROC (NLI clustering), by generator.
  Detection quality decreases as the generator gets larger.}
\label{tab:factualbio}
\input{includes/generated_factualbio_main.tex}
\end{table}

At the claim level the same scale effect from
Section~\ref{sec:bio-results} appears even more cleanly: claim-level
detection AUROC falls monotonically with model size, $0.688$ (Mistral-7B)
$\to 0.644$ (Llama-3.1-8B) $\to 0.521$ (Llama-3.1-70B). By the time we
reach the 70B model, claim-level semantic entropy is barely above chance.
Larger models repeat their confabulations too consistently for a
consistency-based signal to separate true from false claims.

\section{Summary of findings}
\label{sec:results-summary}

The experiments support a single, regime-dependent picture.

\begin{enumerate}
  \item \textbf{Short answers: semantic clustering is close to
    redundant.} All sampling-based detectors work (AUROC
    $\approx 0.74$--$0.88$), but the cheap surface/exact-match baseline
    is the strongest single score in most cells and \emph{no method
    beats it significantly} (Section~\ref{sec:significance}).
    Normalization, not entailment, does most of the grouping.
  \item \textbf{Long-form: semantic clustering is decisive.} Surface
    baselines collapse to chance on biographies, while meaning-aware
    entropy, the LLM-judge variant, and the hidden-state probe all
    recover a real signal (Section~\ref{sec:bio-results}). The clustering
    backend is result-determining here
    (Section~\ref{sec:clustering-sensitivity}).
  \item \textbf{Scale erodes the signal.} Both whole-paragraph and
    claim-level detection get \emph{worse} as the generator grows,
    because larger models confabulate more consistently
    (Section~\ref{sec:factualbio-results}).
  \item \textbf{Temperature decides which detector wins.} Sampling-based
    entropy and the single-pass probe respond to decoding temperature in
    opposite directions and cross over on short answers
    (Section~\ref{sec:temperature-sensitivity}).
  \item \textbf{Two blind spots, seen by hand.} Confident systematic
    errors and noisy gold labels (Example~\ref{sec:audit-artifact}) are
    real limits that the aggregate scores hide; they are taken up in
    Chapter~\ref{ch:discussion}.
\end{enumerate}

\chapter{Discussion}
\label{ch:discussion}

Beating the original baselines is weaker evidence for semantic grouping than
beating normalized-string counting. The original baselines differ from semantic
entropy in several ways at once, so an advantage over them cannot be traced to
any single step. Normalized-string counting differs in one step only, which
makes it the comparison that carries information about grouping. In the
sentence-answer comparisons under LLM grading, the tested Qwen entailment model
clears that stronger comparison, whereas no additional benefit is detected with
the tested DeBERTa models. The entailment model is therefore central to
interpreting the method's performance, and this chapter returns to it in every
section.

\section{What grouping by meaning adds}
\label{sec:what-clustering-is-worth}

\subsection{Normalization gives the most consistent improvement}

The decomposition in Section~\ref{sec:ladder} separates counting strings,
normalizing them, grouping them by meaning, and weighting the resulting groups
by likelihood. Normalization gives the most consistent improvement across the
tested conditions. It adds 0.0398 AUROC on average, it is positive in all 24
configuration entries, and every one of those intervals excludes zero. The two
string scores do not use an entailment model, so the 24 entries repeat 12
distinct condition--grader contrasts, and all 12 are positive. Giving those 12
equal weight still yields $+0.0372$.

Normalization is also the only step whose
difference exceeds the seed-variation reference of 0.0245, which means it is
the only step that moves AUROC further than rerunning the same configuration
under a different sampling seed. What it does is mechanical rather than
semantic: it merges differences in capitalization, punctuation, articles, and
spacing without calling another language model.

The other three steps behave differently. Counting raw strings instead of
reading sequence likelihoods is positive in only 4 of 24 comparisons, and in 8
comparisons the interval favors sequence-likelihood entropy instead. The four
positive differences are large enough to carry the overall mean to $+0.0067$,
which is a case where a positive average does not describe most of the
comparisons behind it. The claim that counting answer strings beats reading
log-likelihoods therefore holds only after the strings are normalized.

Grouping the samples by meaning adds 0.0016 AUROC over normalized-string
counting and is positive in 14 of 24 comparisons, with seven intervals favoring
normalized-string counting. Weighting the meaning groups by likelihood adds
0.0015. Both of these sit far below the normalization step and well inside the
seed range, which explains why discrete semantic entropy performs so similarly
to the likelihood-weighted version throughout this study.

The comparison with surface entropy matters because the original baselines do
not isolate grouping by meaning. Surface entropy also outperforms those
baselines in every comparison where it was measured, so a result reported only
against the original baselines would credit grouping with an improvement that
normalization already delivers. The decomposition identifies normalization,
rather than grouping alone, as the consistent source of improvement.

\subsection{The entailment model decides whether grouping adds anything}

In the short-answer comparisons, no additional benefit over normalized-string
counting is detected with the tested DeBERTa cross-encoders under LLM grading.
None of the corresponding hierarchical intervals excludes zero. This holds
across all twelve comparisons that pair a released cross-encoder with an LLM
correctness judge, so it is not one unlucky configuration.

Substituting \texttt{Qwen2.5-72B-Instruct} as the entailment model changes that
picture. On identical sentence-length answers under the Qwen judge, discrete
semantic entropy records 0.751 with the large DeBERTa-v3 model, 0.757 with the
released-code default, and 0.823 with Qwen-72B, against a surface entropy of
0.7364 that does not move at all because it never calls an entailment model.
The resulting differences are $+0.0147$, $+0.0204$, and $+0.0867$, and only the
last has an interval that excludes zero.

The stored answers are identical in
all three cases, so the entailment model is the only factor that changed. The
spread it produces, from 0.751 to 0.823, is larger than most of the method
differences this thesis measures.

Two checks support reading this as an entailment-model effect rather than an
artifact. Re-grading the same answers with Llama-3.1-70B leaves the ranking
unchanged and the sizes close, at $+0.0872$, $+0.0223$, and $+0.0139$, which
reduces the concern that one model serving as both entailment model and
correctness judge inflates the Qwen result.

Splitting the sentence-length
difference by generator shows the same dependence from another angle. Under the
Llama judge on the default entailment model, the generator with the longest
answers, Llama-3.1-8B at 667 mean characters, has the smallest difference at
$+0.012$. On the Qwen entailment model the same generator has the largest at
$+0.132$. The generations behind those two columns are the same text.

The entailment model is therefore part of the measurement rather than a setting
in the background. The recorded grid varies it across three models, which is
enough to show that the choice moves the result and not enough to say which
choice is closest to human judgments of equivalence. No human equivalence
labels were collected, so the three models can be compared with each other but
not against an outside standard.

The pattern across all 24 comparisons carries the same message. Eleven
hierarchical intervals exclude zero. The four that favor semantic entropy all
use the Qwen entailment model, and the seven that favor surface entropy all use
token-F1 correctness labels. Neither the estimator nor its comparator changes
across those eleven results. Only the entailment model and the correctness rule
do.

\subsection{What answer length explains, and what it does not}
\label{sec:length-discussion}

Farquhar et al.~\cite[Supplementary Note~7 and Supplementary Fig.~7]{semantic_entropy_nature_2024}
report a larger advantage over naive entropy for sentence-length answers than
for short phrases. The comparison in Section~\ref{sec:answer-length-results}
shows the same direction across four prompting conditions. The
discrete-semantic advantage over naive predictive entropy rises from 0.0110 for
short phrases at 11.2 mean characters, to 0.0326 at 69.8 characters, to 0.0628
for the sentence instruction without examples at 255.2 characters.

The ordering
follows the measured answer length rather than the name of the condition, which
is visible in the sentence instruction with five examples: it produces answers
of 11.4 characters and a difference of 0.0094, in line with the other short
condition.

That agreement weakens against the stronger comparator. Measured against
surface entropy instead of naive predictive entropy, the same four differences
run from $-0.0020$ to $+0.0204$, and none of the four hierarchical intervals
excludes zero. Length therefore predicts how much semantic entropy gains over
reading sequence likelihoods, and does not establish a gain over counting
normalized strings with the default entailment model.

The long-form biography evaluation shows the largest recorded gap in the study:
0.670 with large DeBERTa and 0.708 with Qwen, against 0.497 for surface and
naive sample entropy. The reason is visible in the inputs. Nearly every sampled
paragraph is a distinct string, so every string-based group holds one sample,
the frequencies are equal, and the string entropy is constant and near chance
by construction.

Meaning-based grouping restores variation in the score
(Section~\ref{sec:longform-results}), and the pooled gaps are $+0.174$ with
large DeBERTa and $+0.212$ with Qwen. Those intervals resample records while
holding generators and seeds fixed, so they are narrower in scope than the
short-answer intervals. The task, prompts, and grading target also change
together with length, so this result does not isolate answer length.

\section{Why the evaluation setup matters}
\label{sec:apparatus}

The entailment model, sampling temperature, and correctness rule can change both
the size and the direction of a method comparison. They are also different kinds
of choice. The entailment and grading checks reuse the same stored answers, so
the generated text is held fixed and only the analysis moves. The temperature
experiment instead changes the uncertainty samples while keeping the graded
answer and its LLM correctness label fixed. Neither should be described as a
change made to one fixed generation set, because only the first leaves the
samples untouched.

\subsection{The entailment model is part of the method}

Reproducing these comparisons takes three pieces of information: the entailment
model, the prompt, and the clustering rule. The label ``bidirectional
entailment'' alone fixes none of them, and the measurements above show what
that costs. The same answers and the same estimator span 0.751 to 0.823
depending on which model reads them. A reader given only the label cannot tell
where in that range a reported number sits, and cannot tell whether a gap over
a string baseline reflects the estimator or the model chosen to support it.

This thesis records all three for every comparison, which is what makes its
entailment columns comparable to one another. What it does not record is an
outside reference for the clustering rule itself. The rule is the one in the
released code, including its order dependence, and it was kept unchanged so
that the entailment-model comparison would isolate the model, so the study
measures how much the rule's output moves and not how often it is right.

\subsection{Sampling temperature changes which detector leads}
\label{sec:diversity-vs-likelihood}

At lower temperatures the sampled answers become less varied, and the detectors
that measure variation across samples have less to measure. Discrete semantic
entropy falls from 0.757 at $T=1.0$ to 0.568 at $T=0.1$, and surface entropy
from 0.736 to 0.589. Naive predictive entropy, which reads sequence likelihoods
rather than spread across samples, falls much less, from 0.694 to 0.644, and at
$T=0.1$ it leads both diversity-based estimators. The ordering therefore
reverses inside the tested range.

The graded answers and their correctness
labels are identical at every temperature, so the reversal comes from the
uncertainty samples alone. The highest temperature tested here is the most
favorable setting for the sample-based detectors in this sweep rather than a
general optimum, and a comparison reported at one temperature does not carry
over to another.

\subsection{The correctness rule can reverse the comparison}

Identical answers and identical detector scores can produce a different method
ranking when the correctness labels change. In the sentence-length condition,
surface entropy leads under token-F1, while semantic entropy has higher point
estimates under both LLM judges, and the LLM-graded intervals include zero with
the default entailment model. The short-phrase condition shows the same
reversal in miniature. The paired difference is $-0.0133$ under token-F1,
$+0.0005$ under the Qwen judge, and $+0.0027$ under the Llama judge. All seven
hierarchical intervals that favor surface entropy use token-F1 labels.

The mechanism is not subtle. Token-F1 compares an answer with a short
reference, so it can penalize an answer that is correct but long. Its low
acceptance rate in some groups also makes AUROC estimates unstable. Excluding
those groups leaves surface entropy's token-F1 advantage largely unchanged with
the default entailment model (Section~\ref{sec:grading-results}), so the effect
is not confined to a few extreme cases. The grading rule is part of what a
reported AUROC means.

\subsection{The resampling scheme changes which differences survive}

How an interval is built matters as much as what it contains. Across the same
24 comparisons, resampling records alone leaves 18 intervals that exclude zero,
with a mean width of 0.0153. Resampling the dataset--generator units as well
leaves 11, with a mean width of 0.0467, about three times wider. Seven apparent
differences disappear.

Resampling units alone identifies the same eleven, which
places the extra variation between units rather than between questions. With
three datasets and three generators there are only nine units, so a difference
that survives the record-level scheme alone describes the questions in this
grid rather than the grid itself.

\section{Comparison with Farquhar et al.\ and Kossen et al.}
\label{sec:comparison-farquhar}
\label{sec:reproduction}

The sentence-answer results agree with Farquhar et
al.~\cite{semantic_entropy_nature_2024} on the measured original baseline
comparisons and on the similar performance of the two semantic estimators.
Discrete semantic entropy leads P(True) by 0.025 on average over six
comparisons, the semantic entropy probe by 0.051, and the accuracy probe by
0.049, and it leads all three in every comparison where each was measured.
These are the comparisons the original study makes, and they reproduce here.

The agreement is limited to the configurations tested. Farquhar et
al.~\cite{semantic_entropy_nature_2024} report that P(True) tends to improve
with model size. P(True) does not show this trend within the tested Llama
models, where the 70B model records 0.711 against 0.726 for the 8B in
\texttt{chat\_0shot} and 0.819 against 0.826 in \texttt{default\_5shot}. A
comparison within one family and across two sizes is not enough to reject a
broader trend.

The baseline contrasts also do not hold semantic entropy's
entailment model fixed, because the comparator uses the default entailment
model against P(True) and \texttt{nli-deberta-v3-large} against the probes. The
missing saved training feature limits exact reproduction of the probe fits.

Kossen et al.~\cite{kossen_semantic_entropy_probes_2024} evaluate probes both
within their training tasks and on held-out tasks, and their reported advantage
over accuracy probes concerns transfer to unseen tasks. The six probe
comparisons in Section~\ref{sec:paper-baselines} use in-distribution probes,
trained and evaluated on separate questions from the same task. They do not
test that transfer finding, so the probe results here are not evidence for or
against it.

The reimplementation uses three of Farquhar et al.'s five short-answer
datasets: TriviaQA, NQ-Open, and SVAMP. It retains Mistral 7B Instruct and
tests Llama-3.1 at 8B and 70B in place of LLaMA-2 Chat. This gives direct
dataset overlap and partial model overlap while testing a smaller grid. The
generation settings also allow the answer being graded to be longer than the
samples used to estimate uncertainty, so the graded text and the text behind
the uncertainty score are not always produced under the same limits. Following
the released code makes those choices traceable, but does not remove their
effect on the results.

\chapter{Conclusion}
\label{ch:conclusion}

\section{Summary}

This thesis asked whether semantic entropy provides a better
hallucination-risk and abstention signal than surface-form uncertainty,
and how that answer depends on the surface-form variability of the
task, the clustering backend, and the scale of the generator. The
method follows Farquhar et al.~\cite{semantic_entropy_nature_2024};
the contribution here is a reproducible reimplementation evaluated on
a full cross-seed grid --- three open generators, three short-answer
datasets and a long-form biography task, three decoding seeds under
the original paper's sampling and labelling protocol --- plus a
claim-level long-form extension, rather than a new semantic-entropy
method.

The answer is regime-dependent, in both directions and more sharply
than expected. On short-answer tasks every sampling-based estimator
detects errors well (AUROC 0.74--0.88), but entropy over normalized
strings --- the cheapest score in the family --- is the strongest
single method in six of nine model--dataset cells, and no semantic,
kernel, probe, or self-evaluation method beats it significantly
anywhere. A record-level analysis explains why: dataset normalization
already removes almost all surface variation before any semantic step
runs, so clustering has little left to merge. On long-form biography the
same surface scores collapse to chance while every detector that escapes
raw string-counting recovers a statistically significant signal for
every generator; the strongest long-form scores come from
probability-weighted naive entropy on the two smaller models
(AUROC 0.77--0.78) and from the single-pass P(True) on the largest
(0.78). Under the paper-faithful protocol, the short-answer
semantic-entropy averages land in the same band as the original work's
reported 0.790, with the same method ordering.

Two further findings qualify the method's practical use. First, the
clustering backend is result-determining on long-form inputs: across
NLI cross-encoders and an LLM judge, backends span chance-level to
clearly useful behavior on the same saved samples, and the same
artifact reappears at claim level, where exact-match clustering wrongly
makes the largest model look worse than chance. Second, the claim-level FactualBio extension shows that
semantic entropy's discriminative power decreases monotonically with
generator scale (AUROC 0.688 $\rightarrow$ 0.644 $\rightarrow$ 0.521
from 7B to 70B): larger models repeat their confabulations more
consistently, eroding exactly the inconsistency that the method
measures. Hidden-state probes invert this picture. On short answers
they are merely a cheap alternative --- strictly weaker than
sampling-based entropy, and converging to the same performance whether
trained on correctness or on the semantic-entropy threshold --- but on
long-form biography the SEP probe is competitive and, unlike every
sampling-based score, its accuracy \emph{rises} with model scale ---
though naive entropy, and on the 70B model P(True), score higher. Where
a larger model's growing self-consistency erodes entropy- and
agreement-based detection, a probe that reads the model's own
representation does not, making it the natural complement to semantic
entropy precisely where entropy is weakest.

\section{Contributions}

The substantive method comes from Farquhar et
al.~\cite{semantic_entropy_nature_2024}. The thesis contributes six
main artifacts:

\begin{enumerate}
  \item a reproducible \texttt{semantic\_entropy} implementation for
        normalization, clustering under multiple entailment backends,
        the four-estimator entropy family, correctness labeling, and
        selective-prediction metrics;
  \item a complete cross-seed experiment record --- three generators,
        four datasets, three seeds under paper-faithful sampling and
        accuracy labelling --- with Slurm templates, JSONL outputs,
        per-run and cross-seed metric summaries, and record-level
        paired-bootstrap significance kept separate from seed
        stability;
  \item extension baselines covering P(True), in- and
        out-of-distribution embedding regression, Kernel Language
        Entropy under two encoders, and the Semantic Entropy
        Probe~\cite{kossen_semantic_entropy_probes_2024}, with the
        probe-equivalence observation that correctness-trained and
        SE-trained probes are interchangeable in practice;
  \item a long-form evaluation at two granularities: the
        whole-paragraph biography task under the same
        baseline-versus-semantic design, and the claim-level
        FactualBio extension --- biographies generated by the models
        under test, decomposed into claims, re-queried, and scored by
        context-conditioned NLI-clustered semantic entropy with
        judge-graded labels --- which yields the scale trend and the
        clustering-artifact result;
  \item record-level worked audits drawn from the canonical runs that
        trace where clustering actually happens (normalization versus
        NLI), show a confident systematic error the method cannot catch,
        and explain why the surface baseline collapses on long text,
        providing mechanistic rather than purely metric-level evidence
        for the main conclusions (Section~\ref{sec:hand-audits});
  \item a sampling-temperature sweep ($T\in\{0.1,0.3,0.5,0.7\}$ against
        the canonical $T=1.0$) over the full short-answer and biography
        grid, with the numbers recomputed from raw scores as an
        independent check, isolating the opposite temperature dependence
        of the two detector families --- sampling-based entropy
        improving as the decoder warms, the hidden-state probe declining
        --- and the regime-dependent crossover this implies for
        deployment (Section~\ref{sec:temperature-sensitivity}).
\end{enumerate}

The full source code, result tables, Slurm templates, and LaTeX thesis
source are available at \url{https://github.com/D3prave/thesis}.
The Python implementation lives in the \texttt{src/semantic\_entropy/}
package; entry points are documented in \texttt{README.md}.

\section{Limitations}

The thesis is deliberately narrow in scope. The detection target is the
confabulation setting defined in Chapter~\ref{ch:background}.
High-confidence systematic errors are by construction not detected by
this method --- a limitation the FactualBio scale trend makes concrete,
since consistent confabulation at 70B is precisely the confident-error
case; Simhi et al.~\cite{simhi_choke_2025} discuss this setting.

Within the empirical results, the main residual limitations are the
automated LLM judging behind every long-form label (whole-paragraph
and claim-level alike), the three-seed basis of the stability
estimates, and the English-only, biography-only long-form scope.
Record-level issues --- alias-type oracle failures on TriviaQA,
semantically consistent refusals on biography, and SVAMP's
integer-arithmetic structure --- are documented in the audits; the
fuller discussion is in Section~\ref{sec:discussion-limitations}.

\section{Future Work}

Three extensions follow directly from the limitations above. First,
the automated claim and paragraph labels should be validated against a
human-graded
subsample, since every long-form conclusion inherits the judge's error
modes. Second, the FactualBio scale trend invites extension beyond 70B
and beyond biographical claims: if consistency-based detection
continues to erode with scale, the practical future of confabulation
detection lies with methods that do not require the model to re-roll
its errors, such as hidden-state probes or retrieval-grounded
verification. Third, the long-form evaluation should extend beyond
English single-reference biography and be paired with a dedicated
detector for the semantically consistent refusals that low entropy
cannot flag.


\appendix
\chapter{Reproducibility Details}
\label{app:method-details}

\section{Code and data files}
\label{app:schema}
\label{app:reproducibility}

The accompanying repository is at \url{https://github.com/D3prave/thesis}.

Most result tables and figures were generated from run-level measurements
for the short-answer experiment and the biography comparison in
Section~\ref{sec:longform-results}. Each measurement identifies one prompting
condition, dataset, generator, seed, temperature, entailment model,
correctness rule, and uncertainty method. Additional summaries describe the
variation across seeds and the agreement between the two LLM judges. The appendix tables
summarize these run values.

The hierarchical intervals of the short-answer comparisons come from a
record-level paired bootstrap analysis. Rebuilding tables from the resulting
summaries does not recompute the values from the raw answers.

Generation and LLM judging ran on the NHR@FAU Alex GPU cluster. Cross-encoder
clustering ran on Alex and on the Fritz CPU cluster. Given stored answers,
cross-encoder clustering and exact matching are deterministic. LLM entailment
decisions can vary with inference batching.

\section{Long-form biography records}
\label{app:longform-details}

The biography evaluation contains three generators, three seeds, and
500 subjects per generator--seed run. Prompts follow the FActScore subject
format~\cite{min_factscore_2023}. The system instruction requests a factual
biography of about 120--150 words, and the query is ``Tell me a bio of
\{entity\}.'' Ten responses are sampled using vLLM at $T=1.0$, top-$p=0.9$,
top-$k=50$, with at most 256 new tokens each. The saved Wikipedia introduction
provides the reference.

One paragraph per prompt is graded: the most frequent exact normalized
surface form, with ties resolved by first occurrence. Qwen2.5-72B-Instruct
returns one \texttt{CORRECT}/\texttt{INCORRECT} verdict for that paragraph.
Specific facts must be supported by the reference, while vague text without a
contradiction is permitted. The reported acceptance rate is the fraction of
selected paragraphs receiving \texttt{CORRECT}.

The recorded comparisons cluster normalized paragraphs into connected
components. DeBERTa uses answer-only input and requires entailment in both
directions. Qwen receives the question and asks whether both paragraphs convey
the same factual claims, requiring agreement in both response orders. This is
the biography pipeline's grouping rule.
The short-answer pipeline uses the raw-answer anchor scan in
Section~\ref{sec:clustering}. The four settings are exact matching,
DeBERTa-v3-base~\cite{sentence_transformers_nli_deberta_v3_base_2022},
DeBERTa-v3-large, and Qwen-72B. Sequence log-probabilities were not retained
because the vLLM path exposed tempered, truncated probabilities, so this
comparison uses count-based entropy only. Figure~\ref{fig:longform-comparison}
and Tables~\ref{tab:longform}--\ref{tab:longform-by-generator} report the results.

The records-only intervals were computed from the saved labeled biography
records using the same AUROC and paired bootstrap functions as the
short-answer analysis.

For each of 2{,}000 resamples, with random seed 20260831, the analysis draws
500 biographies with replacement within each generator--seed run. Both
estimators and the correctness label use the same sampled indices. Draws are
independent across runs. It averages the paired AUROC differences over the
three seeds, then equally over the three generators for the pooled result.
The 95\,\% intervals are the 2.5th and 97.5th percentiles of these differences.
All 108 relevant run-level AUROCs match the recorded measurements to six decimals.
No resample had an undefined AUROC.

The dataset, generators, and seeds remain fixed. With one dataset and three
dataset--generator units, no unit or hierarchical intervals were computed.
These are the permissive records-only intervals described in
Section~\ref{sec:statistics}, not estimates of uncertainty across
new datasets and generators. The evaluation also uses one correctness judge,
Qwen, which is one of the entailment models.

\section{Collection checks}
\label{app:collection-checks}

\paragraph{Temperature-sweep labels.}
LLM correctness labels are matched by prompt identifier to the main
collection. They are bit-identical across all 108 checked cells, so all sweep
temperatures are evaluated against the same temperature-$0.1$ answer.
Token-F1 instead grades the modal sampled answer because the sweep did not
save a separate most-likely answer. Its accuracy differs from the main run in
all 54 checked cells, by 0.015 on average in absolute terms and at most 0.262.
The reported sweep therefore uses only LLM correctness labels, and its
answer-length fields are not compared across temperatures.

\paragraph{TriviaQA deduplication.}
\label{app:dedup}
The SQuAD-format conversion stores multiple question--context rows for some
questions, and its row-level split places copies of 16 questions in both
training and evaluation. A further nine questions repeat within training and
six within evaluation. Removing duplicate questions leaves 391 training and
378 evaluation questions. The same evaluation exclusions apply to every
method, entailment model, grader, and record-level resampling analysis.
NQ-Open and SVAMP contain no repeats in the selected sets.

\section{Bootstrap details}
\label{sec:bootstrap-assumptions}

The saved short-answer paired analysis uses 2,000 resamples. Within each selected
dataset--generator pair, it averages the three seed-level AUROC differences,
then averages over the nine selected pairs. Record resampling keeps the
compared methods paired within a seed cell, but draws question indices
separately across seeds and generators. It neither synchronizes repeated
questions across these cells nor resamples datasets and generators as
separate factors.

A record resample with only one correctness class has undefined AUROC and is
omitted from that draw's average. This can matter when few answers are labeled
correct. Omission rates were not retained in the original export. The additional
analyses in Appendix~\ref{app:robustness} report them explicitly. The original
intervals and bootstrap tail
proportions are conditional on these choices and are not adjusted for
multiple comparisons. The released code instead reports 90\,\% bias-corrected
and accelerated intervals from 1{,}000 resamples, so its intervals are not
directly comparable with these.

Run-level measurements can verify paired point estimates, but do not suffice
to recompute the record-level intervals, which require the scored answers.
Appendix~\ref{app:robustness} reports additional record-level
analyses.

\section{Probe training and saved features}
\label{app:probe-training}

For each training question, the most-likely answer supplies a hidden state.
Ten separate sampled responses supply the discrete semantic entropy target.
The high/low threshold minimizes the sum of squared deviations within the
two entropy groups, following Equation~(5) of
Kossen et al.~\cite{kossen_semantic_entropy_probes_2024}. A logistic-regression
classifier predicts that label from the hidden state. The accuracy probe
instead learns correctness labels. At evaluation time, each probe needs one
answer and its hidden state.

Training is separate for each dataset, generator, and seed, using questions
disjoint from evaluation. Feature choices are fixed rather than tuned on the
evaluation set. The hidden state comes from position
$n_{\mathrm{generated}}-1$, where the stop sequence determines
$n_{\mathrm{generated}}$. An end-of-sequence token need not be present.

All 108 probes use final-layer features, last-token pooling,
$L_2$ regularization with $C=1$, a training seed of zero and at most 2{,}000
optimizer iterations, fitted with the \texttt{lbfgs} solver without feature
standardization. All fits report convergence. Feature dimensions are 4{,}096
for the two smaller models and 8{,}192 for Llama-3.1-70B.

The accuracy probe learns the stored training correctness labels. Scoring it
with another judge does not retrain it. The stored training records do not
say which judge produced these labels.

For the NQ-Open Llama-3.1-70B short-phrase run with seed 42, both probes report
399 training examples and one excluded prompt, which matches the one training
record without a saved feature. All evaluation records have saved features.

\chapter{Secondary and Exploratory Results}
\label{app:secondary-results}

This appendix gives the complete results behind the secondary and exploratory
analyses summarized in Chapter~\ref{ch:results}. The same aggregation script
used for the main results assembles the tables. Point estimates come from
\texttt{measurements.csv}. The hierarchical confidence intervals in
Table~\ref{tab:all-slices} are imported from the record-level paired-bootstrap
analysis because they cannot be recomputed from the exported measurement table
alone.

As in the main Results chapter, a comparison is one reported combination of
prompting condition, entailment model, correctness rule, and pair of methods.
An AUROC difference is the first method's AUROC minus the second's, so a
positive value favors the first method.

\section{All semantic-versus-surface comparisons}
\label{app:all-comparisons}

Table~\ref{tab:all-slices} and Figure~\ref{fig:comparison-forest} report all 24
comparisons. Eleven hierarchical intervals exclude zero. Four favor semantic
entropy, and all four use the Qwen entailment model. Seven favor surface
entropy, and all seven use token-F1 correctness labels.

The two smaller Qwen-entailment-model differences that favor semantic entropy are below
the 0.025 variation observed across generation seeds. This is not a
contradiction. The seed range describes how much one method's AUROC changes
across seeds. A paired test compares two methods on the same generated answers,
so variation shared by both methods cancels when their scores are subtracted.
The seed range is therefore useful for judging one AUROC value, while the paired
interval is the correct check for a difference between methods.

% Generated by scripts/build_thesis_tables.py from docs/independent_analysis/data/measurements.csv. Do not edit by hand.
\begin{table}[htb]
\thisfloatsetup{capposition=above}
\centering
\footnotesize
\caption{All 24 comparisons: discrete semantic entropy minus surface entropy, with hierarchical intervals resampling both the nine units and the records within them. The final column names the direction where the interval excludes zero. Every comparison favoring semantic entropy uses the Qwen entailment model; every comparison favoring surface entropy uses token-F1 grading.}
\label{tab:all-slices}
\begingroup
\setlength{\tabcolsep}{3pt}
\begin{tabular}{@{}lllccl@{}}
\toprule
Condition & \makecell{Entailment\\model} & Grader & \makecell{AUROC\\difference} & \makecell{95\,\% CI\\(hierarchical)} & \makecell{Supported\\direction} \\
\midrule
\texttt{chat\_0shot} & \texttt{xlarge} & Qwen judge & +0.0204 & [-0.030, +0.059] & --- \\
\texttt{chat\_0shot} & \texttt{xlarge} & Llama judge & +0.0223 & [-0.033, +0.066] & --- \\
\texttt{chat\_0shot} & \texttt{xlarge} & token-F1 & -0.0506 & [-0.084, -0.022] & surface \\
\texttt{chat\_0shot} & \texttt{v3-large} & Qwen judge & +0.0147 & [-0.027, +0.050] & --- \\
\texttt{chat\_0shot} & \texttt{v3-large} & Llama judge & +0.0139 & [-0.039, +0.055] & --- \\
\texttt{chat\_0shot} & \texttt{v3-large} & token-F1 & -0.0587 & [-0.091, -0.029] & surface \\
\texttt{chat\_0shot} & \texttt{Qwen-72B} & Qwen judge & +0.0867 & [+0.060, +0.119] & SE \\
\texttt{chat\_0shot} & \texttt{Qwen-72B} & Llama judge & +0.0872 & [+0.056, +0.124] & SE \\
\texttt{chat\_0shot} & \texttt{Qwen-72B} & token-F1 & -0.0658 & [-0.144, -0.004] & surface \\
\texttt{default\_0shot} & \texttt{xlarge} & Qwen judge & +0.0032 & [-0.026, +0.028] & --- \\
\texttt{default\_0shot} & \texttt{xlarge} & Llama judge & +0.0084 & [-0.021, +0.033] & --- \\
\texttt{default\_0shot} & \texttt{xlarge} & token-F1 & -0.0263 & [-0.045, -0.012] & surface \\
\texttt{chat\_5shot} & \texttt{xlarge} & Qwen judge & -0.0020 & [-0.014, +0.008] & --- \\
\texttt{chat\_5shot} & \texttt{xlarge} & Llama judge & -0.0005 & [-0.011, +0.010] & --- \\
\texttt{chat\_5shot} & \texttt{xlarge} & token-F1 & -0.0174 & [-0.027, -0.007] & surface \\
\texttt{default\_5shot} & \texttt{xlarge} & Qwen judge & +0.0005 & [-0.007, +0.008] & --- \\
\texttt{default\_5shot} & \texttt{xlarge} & Llama judge & +0.0027 & [-0.006, +0.012] & --- \\
\texttt{default\_5shot} & \texttt{xlarge} & token-F1 & -0.0133 & [-0.022, -0.005] & surface \\
\texttt{default\_5shot} & \texttt{v3-large} & Qwen judge & +0.0014 & [-0.007, +0.009] & --- \\
\texttt{default\_5shot} & \texttt{v3-large} & Llama judge & +0.0047 & [-0.005, +0.015] & --- \\
\texttt{default\_5shot} & \texttt{v3-large} & token-F1 & -0.0130 & [-0.023, -0.004] & surface \\
\texttt{default\_5shot} & \texttt{Qwen-72B} & Qwen judge & +0.0119 & [+0.003, +0.021] & SE \\
\texttt{default\_5shot} & \texttt{Qwen-72B} & Llama judge & +0.0139 & [+0.003, +0.024] & SE \\
\texttt{default\_5shot} & \texttt{Qwen-72B} & token-F1 & -0.0064 & [-0.017, +0.004] & --- \\
\bottomrule
\end{tabular}
\endgroup
\end{table}


\begin{figure}[htb]
\centering
\includegraphics[width=\linewidth]{margin_forest.pdf}
\caption{All 24 semantic-versus-surface comparisons. The hierarchical
  intervals resample both the nine (dataset, model) units and the records within
  them. Filled markers identify the eleven intervals that exclude zero. The
  twelve comparisons that use a released cross-encoder with an LLM correctness
  judge appear at the bottom; none of their intervals excludes zero.}
\label{fig:comparison-forest}
\end{figure}

One token-F1 comparison does not exclude zero:
\texttt{default\_5shot} on the Qwen entailment model, with an interval of
$[-0.017, +0.004]$. This condition has the shortest answers and the strongest
entailment model, leaving little room for the two estimators to differ.

Every point estimate in Table~\ref{tab:all-slices} agrees to four decimal
places with the cell-level values behind Table~\ref{tab:sentence-length}. The
two tables are produced through different code paths, so this agreement checks
both paths.

\section{Secondary prompting-factorial analysis}
\label{sec:factorial-results}

Table~\ref{tab:factorial} reports the four prompting combinations on the
default entailment model, averaged over the two LLM correctness judges. The AUROC
difference between semantic entropy and surface entropy is $+0.021$ at 255.2
mean characters, $+0.006$ at 69.8, $-0.001$ at 11.4, and $+0.002$ at 11.2.
Every interval includes zero, and the largest difference is below the
seed-variation range.

% Generated by scripts/build_thesis_tables.py from docs/independent_analysis/data/measurements.csv. Do not edit by hand.
\begin{table}[htb]
\thisfloatsetup{capposition=above}
\centering
\footnotesize
\caption{The $2\times2$ on the canonical backend, mean of the two LLM judges. The isolating contrast is the pair of 0-shot corners, where demonstrations are held at zero and only the instruction changes. Every interval spans zero.}
\label{tab:factorial}
\begingroup
\setlength{\tabcolsep}{3pt}
\begin{tabular}{@{}lccrc@{}}
\toprule
Condition & Mean chars & Demos & SE $-$ surface & 95\,\% CI (units) \\
\midrule
\texttt{chat\_0shot} & 255.2 & 0 & +0.021 & [-0.029, +0.063] \\
\texttt{default\_0shot} & 69.8 & 0 & +0.006 & [-0.022, +0.031] \\
\texttt{chat\_5shot} & 11.4 & 5 & -0.001 & [-0.012, +0.008] \\
\texttt{default\_5shot} & 11.2 & 5 & +0.002 & [-0.006, +0.009] \\
\bottomrule
\end{tabular}
\endgroup
\end{table}


The two zero-shot conditions isolate the instruction because both omit worked
examples. Changing only the instruction changes mean answer length by a factor
of 3.7, while the AUROC difference falls from $+0.021$ to $+0.006$. In the two
five-shot conditions, changing only the instruction changes mean answer length
by 0.2 characters.

Two limits from Section~\ref{sec:conditions} make this a secondary analysis.
Both short-answer conditions are five-shot, so the design separates length from
example count only among the longer answers. The added conditions also exist
only on the default entailment model, so the factorial cannot test the entailment model with
the largest AUROC differences in Section~\ref{sec:clustering-sensitivity}.

\section{Exploratory long-form experiment}
\label{sec:longform-results}

This experiment is exploratory. It covers one biography task and one
correctness judge. The recorded results contain 500 biographies in each of nine
model--seed runs, but Section~\ref{sec:longform-design} explains that the exact entity list
cannot be recovered from the stored repository. The default entailment model was not
run, and this study uses a stricter paragraph-clustering rule than the original
study.

Table~\ref{tab:longform} reports those runs under claim-level grading by the
Qwen judge. Mean claim accuracy is 0.375.

\begin{table}[htb]
\thisfloatsetup{capposition=above}
\centering
\footnotesize
\caption{The long-form arm: nine cells, 500 biographies each, claim-level grading by the Qwen judge. Strict bidirectional entailment throughout, which is not the rule the original study applies to paragraphs (Section~\ref{sec:longform-design}).}
\label{tab:longform}
\begingroup
\setlength{\tabcolsep}{3pt}
\begin{tabular}{@{}lccc@{}}
\toprule
Backend & Discrete SE & Surface & Naive sample \\
\midrule
\texttt{exact-match} & 0.4965 & 0.4965 & 0.4973 \\
\texttt{nli-deberta-v3-base} & 0.5245 & 0.4965 & 0.4973 \\
\texttt{nli-deberta-v3-large} & 0.6703 & 0.4965 & 0.4973 \\
\textbf{\texttt{Qwen2.5-72B-Instruct}} & 0.7084 & 0.4965 & 0.4973 \\
\bottomrule
\end{tabular}
\endgroup
\par\vspace{2pt}\footnotesize Mean claim accuracy 0.375 over the nine cells. Surface and naive sample entropy are constant across backends because they never consult one; they are not performing badly but are structurally degenerate, since ten sampled biographies are never string-identical and every sample forms its own cluster. The canonical \texttt{deberta-v2-xlarge-mnli} was not applied to this arm.
\end{table}


Surface entropy scores 0.4965 and naive sample entropy 0.4973 under every
entailment model. Both are near chance. These baselines have a structural limit
for varied paragraph outputs: whenever all ten samples are different strings,
each forms its own group and the entropy reaches the same maximum value. The
aggregate AUROC values do not show that every record has a constant score.
In particular, a fully constant score would give AUROC exactly 0.5 with
average-rank handling of ties. These near-chance baselines provide only a weak
test of the semantic estimator.

The results also show that the choice of entailment model matters.
Discrete semantic entropy has an AUROC of 0.5245 with the base DeBERTa-v3 model,
0.6703 with the large variant, and 0.7084 with Qwen-72B.
This 0.184 range shows strong dependence on the entailment model,
but it is not a clean ranking of paragraph-entailment models. The experiment
uses \texttt{strict\_entailment} throughout, while the original paragraph
experiment uses the relaxed rule. A strict rule with a lower-recall model can
split answers too aggressively. The single Qwen correctness judge is also used
for clustering, so this experiment lacks a check with a different grader.
The original study reports its strongest paragraph result with DeBERTa, but
the different clustering rules mean that the current result is not a direct
test of that ranking.

\section{Secondary comparison of AUROC and AURAC}
\label{sec:aurac-disagreement}

AUROC is the primary metric in this thesis. Table~\ref{tab:aurac} provides a
secondary sensitivity analysis using the two rejection-accuracy definitions.
With AUROC, discrete semantic entropy leads surface entropy in 14 of 24
comparisons, with a mean difference of $+0.0016$. With the released
\texttt{aurac\_paper} definition, it leads in 8 of 24 and the mean difference
is $-0.0015$. With \texttt{aurac\_mean\_retained}, it leads in 6 of 24 and the
mean difference is $-0.0033$.

\input{includes/t_aurac.tex}

Across the 648 discrete-semantic-entropy cells at $T=1.0$,
\texttt{aurac\_paper} correlates with the generating model's accuracy at
$r = 0.971$, and \texttt{aurac\_mean\_retained} correlates at $r = 0.965$.
The corresponding correlation for AUROC is $r = 0.502$. The two AURAC measures
therefore reflect the model's accuracy much more strongly than AUROC does.


\backmatter

\nocite{semantic_entropy_nature_2024}
\printbibliography[heading=bibintoc]

\end{document}
